% 线性代数 — 中文教材草稿.
% Engine: XeTeX (tectonic / xelatex). Chinese via ctexbook with the fandol font set,
% which renders identically on macOS, Linux, and CI.
\documentclass[11pt, fontset=fandol]{ctexbook}

\usepackage{amsmath, amssymb, amsthm, mathtools}
\usepackage{extarrows}   % \xlongequal{...}
\usepackage{mathdots}    % \iddots
\usepackage{tikz}
\usetikzlibrary{arrows.meta, calc}
\usepackage[a4paper, top=2.6cm, bottom=2.6cm, left=2.4cm, right=2.4cm]{geometry}
\usepackage{enumitem}
\usepackage{xcolor}      % color mixing (blue!45!black) for hyperref links
\usepackage{tcolorbox}
\tcbuselibrary{skins, breakable}
\usepackage[colorlinks=true, linkcolor=blue!45!black, citecolor=blue!45!black, urlcolor=blue!45!black]{hyperref}
\hypersetup{bookmarksnumbered=true, pdfstartview=FitH}

% ---------- theorem-like environments ----------
\theoremstyle{plain}
\newtheorem{theorem}{定理}[chapter]
\newtheorem{lemma}[theorem]{引理}
\newtheorem{corollary}[theorem]{推论}
\newtheorem{proposition}[theorem]{命题}
\newtheorem{example}{例}[chapter]
\newtheorem{property}{性质}[chapter]

\theoremstyle{definition}
\newtheorem{definition}{定义}[chapter]

\theoremstyle{remark}
\newtheorem{remark}{注}[chapter]

\renewcommand{\proofname}{证}

% ---------- macros ----------
\newcommand{\trans}{^{\mathrm{T}}}                 % transpose: A\trans
\newcommand{\R}{\mathbb{R}}
\DeclareMathOperator{\rank}{r}                     % \rank(A) prints r(A)

% End-of-chapter exercise heading; #1 is the Chinese chapter numeral.
\newcommand{\chapterexercises}[1]{%
  \section*{习题#1}
  \addcontentsline{toc}{section}{习题#1}
}

\setlength{\parskip}{2pt}

% Question box: the title is a question readers commonly stumble on; the body answers it.
\newtcolorbox{qa}[1]{enhanced, breakable,
  colback=gray!7, colframe=black!60, boxrule=0.5pt, arc=2pt,
  left=7pt, right=7pt, top=5pt, bottom=5pt,
  title={问：#1}, coltitle=white, colbacktitle=black!75,
  fonttitle=\bfseries\small}

% Insight box for the geometric / essential viewpoint.
\newtcolorbox{insight}[1][本质]{enhanced, breakable,
  colback=yellow!10, colframe=yellow!55!black, boxrule=0.5pt, arc=2pt,
  left=7pt, right=7pt, top=5pt, bottom=5pt,
  title={#1}, coltitle=white, colbacktitle=yellow!45!black,
  fonttitle=\bfseries\small}

\title{\huge 线性代数}
\author{《线性代数》编写组}   % TODO: replace with the real author name
\date{2026年10月 \quad 草稿 v0.1（供审阅）}

\begin{document}

\maketitle

\frontmatter
\chapter*{前言}
\addcontentsline{toc}{chapter}{前言}

线性代数是高等学校理工科各专业的一门重要基础课，它的研究对象是向量、矩阵、
线性方程组、线性变换与二次型等。线性代数不仅是学习后续课程（如概率统计、
数值分析、运筹学、机器学习等）的必备工具，其公理化与结构化的思想方法本身
也是数学思维训练的重要内容。

本书按照标准的教学大纲编写，全书共六章：

\begin{itemize}[itemsep=1pt]
  \item 第一章\quad 行列式——行列式的定义、性质、展开定理与克拉默法则；
  \item 第二章\quad 矩阵及其运算——矩阵运算、逆矩阵、分块矩阵、初等变换与矩阵的秩；
  \item 第三章\quad 向量组的线性相关性——线性组合、线性相关性、向量组的秩与向量空间；
  \item 第四章\quad 线性方程组——解的判定定理、基础解系与解的结构；
  \item 第五章\quad 特征值与特征向量——内积与正交性、方阵的对角化、实对称矩阵的正交对角化；
  \item 第六章\quad 二次型——二次型的标准形、惯性定理与正定二次型。
\end{itemize}

每章末配有习题，书末附录给出了部分习题的答案或提示。本书只假定读者具备高中
数学知识；第一章从二元、三元线性方程组出发引入行列式，各章内容循序渐进，
教师可根据学时适当取舍（例如第三章与第四章的顺序可以互换讲授）。

本先行版包含第一至第五章，并配有“问与答”框与几何直观插图，以帮助读者理解
概念的本质；第六章与习题答案附录将在后续版本中陆续补齐。

本书为草稿（v0.1），全部内容开放审阅。读者若发现概念表述、公式、例题演算或
习题答案中的任何问题，欢迎通过代码托管平台的问题跟踪（Issues）页面逐条指出，
并请注明章节与题号，以便修订。

\begin{flushright}
编者\quad 2026年10月
\end{flushright}

\tableofcontents

\mainmatter
\chapter{行列式}

行列式是线性代数中最基本的概念之一，它起源于求解线性方程组的需要。本章从
二元、三元线性方程组的求解公式引出二阶、三阶行列式，进而推广到 $n$ 阶行列式，
讨论其性质与计算方法，最后介绍用行列式求解线性方程组的克拉默法则。

\section{二阶与三阶行列式}

考虑二元线性方程组
\[
  \begin{cases}
    a_{11}x_1 + a_{12}x_2 = b_1, \\
    a_{21}x_1 + a_{22}x_2 = b_2.
  \end{cases}
\]
用消元法可得：当 $a_{11}a_{22}-a_{12}a_{21}\neq 0$ 时，
\[
  x_1=\frac{b_1a_{22}-a_{12}b_2}{a_{11}a_{22}-a_{12}a_{21}},
  \qquad
  x_2=\frac{a_{11}b_2-b_1a_{21}}{a_{11}a_{22}-a_{12}a_{21}} .
\]
两个分母恰好是同一个代数和。为了简明地表达这类结果，给出下述定义。

\begin{definition}\label{def:order2}
称代数和 $a_{11}a_{22}-a_{12}a_{21}$ 为二阶行列式，记作
\[
  \begin{vmatrix} a_{11} & a_{12} \\ a_{21} & a_{22} \end{vmatrix}
  = a_{11}a_{22}-a_{12}a_{21},
\]
其中 $a_{ij}$（$i=1,2$ 为行标，$j=1,2$ 为列标）称为行列式的\textbf{元素}。
\end{definition}

二阶行列式可用“对角线法则”记忆：主对角线（左上至右下）两元素之积取正号，
副对角线（右上至左下）两元素之积取负号。

\begin{qa}{行列式到底是一个数，还是一张表？为什么只有方阵才能谈行列式？}
答：行列式是一个\textbf{数}：它把 $n^2$ 个数（一张方表）按确定规则压缩成
一个数。因此 $|A|$ 能参加普通四则运算，而 $A$ 本身不能。非方阵没有行列式，
因为定义要求“每行每列各取一个元素、恰好取尽 $n$ 个”，行列数不同时无法
配对；从稍后的几何观点看，“面积（体积）缩放因子”只对维数相同的空间之间
的变换才有意义。
\end{qa}

\begin{example}
计算二阶行列式 $\begin{vmatrix} 3 & 1 \\ 2 & 4 \end{vmatrix}$。
\end{example}
\begin{proof}[解]
$\begin{vmatrix} 3 & 1 \\ 2 & 4 \end{vmatrix}=3\times 4-1\times 2=10$。
\end{proof}

类似地，解三元线性方程组时会出现形如
\[
  a_{11}a_{22}a_{33}+a_{12}a_{23}a_{31}+a_{13}a_{21}a_{32}
  -a_{13}a_{22}a_{31}-a_{11}a_{23}a_{32}-a_{12}a_{21}a_{33}
\]
的六项代数和，称其为\textbf{三阶行列式}，记作
\[
  \begin{vmatrix}
    a_{11} & a_{12} & a_{13} \\
    a_{21} & a_{22} & a_{23} \\
    a_{31} & a_{32} & a_{33}
  \end{vmatrix}.
\]
三阶行列式也可按对角线法则展开：三条主对角线方向（平行于主对角线）之积取
正号，三条副对角线方向之积取负号。注意对角线法则只适用于二阶、三阶行列式，
更高阶的行列式需要一般的定义（见 \S\ref{sec:n-order}）。

\begin{example}
计算三阶行列式
$D=\begin{vmatrix} 1 & 2 & 3 \\ 4 & 5 & 6 \\ 7 & 8 & 10 \end{vmatrix}$。
\end{example}
\begin{proof}[解]
\begin{align*}
  D &= 1\cdot 5\cdot 10 + 2\cdot 6\cdot 7 + 3\cdot 4\cdot 8
       - 3\cdot 5\cdot 7 - 1\cdot 6\cdot 8 - 2\cdot 4\cdot 10 \\
    &= (50+84+96)-(105+48+80) = 230-233 = -3 .
\end{align*}
\end{proof}

\subsection{几何直观：行列式是面积的缩放因子}

把二阶行列式的两个列向量 $(a_{11},a_{21})\trans$ 与 $(a_{12},a_{22})\trans$
画在平面上，它们张成一个平行四边形。把这个平行四边形沿对角线切开拼补，
或直接用顶点坐标计算三角形面积，可以验证其面积恰为
$\Bigl|\begin{vmatrix} a_{11} & a_{12} \\ a_{21} & a_{22} \end{vmatrix}\Bigr|$。
于是：行列式的\textbf{绝对值}是“面积缩放因子”，\textbf{符号}表示“定向”——
第一列转到第二列为逆时针时取正，顺时针时取负。

\begin{figure}[htbp]
\centering
\begin{tikzpicture}[>=Stealth, scale=1.05]
  \begin{scope}
    \draw[->] (-0.4,0) -- (2.4,0) node[below] {$x_1$};
    \draw[->] (0,-0.4) -- (0,2.4) node[left] {$x_2$};
    \fill[blue!12] (0,0) -- (1,0) -- (1,1) -- (0,1) -- cycle;
    \draw[very thick, blue!60!black, ->] (0,0) -- (1,0) node[below right] {$\varepsilon_1$};
    \draw[very thick, blue!60!black, ->] (0,0) -- (0,1) node[left] {$\varepsilon_2$};
    \node[below] at (1,-0.75) {标准基张成的单位正方形，面积 $1$};
  \end{scope}
  \draw[->, very thick, gray] (3.1,1) -- (4.5,1)
        node[midway, above] {$A$};
  \begin{scope}[xshift=5.6cm]
    \draw[->] (-0.4,0) -- (3.2,0) node[below] {$y_1$};
    \draw[->] (0,-0.4) -- (0,2.4) node[left] {$y_2$};
    \fill[red!12] (0,0) -- (2,0) -- (3,1) -- (1,1) -- cycle;
    \draw[very thick, red!70!black, ->] (0,0) -- (2,0) node[below] {$(2,0)\trans$};
    \draw[very thick, red!70!black, ->] (0,0) -- (1,1) node[above left] {$(1,1)\trans$};
    \draw[gray!60, dashed] (2,0) -- (3,1);
    \draw[gray!60, dashed] (1,1) -- (3,1);
    \node[below] at (1.5,-0.75) {列向量张成的平行四边形，面积 $2$};
  \end{scope}
\end{tikzpicture}
\caption{矩阵 $A=\begin{pmatrix} 2 & 1 \\ 0 & 1 \end{pmatrix}$ 把单位正方形
（面积 $1$）变成其列向量张成的平行四边形（面积 $2=|A|$）。行列式就是这次
变换的面积缩放因子，符号记录定向是否翻转。}
\label{fig:det-area}
\end{figure}

用这一幅图可以“看出”本章后面的许多性质：倍加变换（性质~\ref{prop:det-add}）
相当于把平行四边形沿一条边的方向“推斜”，底不变、高不变，面积不变；两行
相同（性质~\ref{prop:det-equal-rows}）相当于两条邻边重合，图形被压扁成线段，
面积为 $0$；某行乘 $k$（性质~\ref{prop:det-factor}）相当于把一条边拉长 $k$
倍，面积随之变为 $k$ 倍。三阶行列式同理是平行六面体的定向体积，$n$ 阶行列式
是 $n$ 维“平行多面体”体积的缩放因子。

\section{$n$ 阶行列式的定义}\label{sec:n-order}

为了把行列式推广到任意阶，需要排列的概念。

\begin{definition}\label{def:permutation}
把 $n$ 个不同的元素排成一列，称为这 $n$ 个元素的一个\textbf{排列}。
$n$ 个不同元素的所有排列共有 $n!$ 个。在一个排列中，若一个较大的数排在
一个较小的数前面，就称这两个数构成一个\textbf{逆序}；一个排列中逆序的总数
称为该排列的\textbf{逆序数}，记作 $t$。逆序数为奇数的排列称为\textbf{奇排列}，
逆序数为偶数的排列称为\textbf{偶排列}。
\end{definition}

例如排列 $3\,2\,1$ 中，逆序有 $(3,2)$、$(3,1)$、$(2,1)$，故 $t=3$，它是奇排列；
自然排列 $1\,2\,3$ 的逆序数为 $0$，是偶排列。

\begin{theorem}\label{thm:swap-parity}
对换一个排列中的两个元素，排列改变奇偶性。
\end{theorem}
\begin{proof}
先看相邻对换：设排列 $\cdots i\,j\cdots$ 对换为 $\cdots j\,i\cdots$。此时只有
$i,j$ 的相互顺序改变，其余元素间的逆序关系不变。若 $i,j$ 原本构成逆序，
对换后逆序数减少 $1$；否则增加 $1$。故相邻对换使逆序数增减 $1$，奇偶性改变。

再看一般对换：设两元素之间相隔 $2k$ 个元素（相距 $2k+1$ 个位置）。先把前一
元素与其后的 $2k$ 个元素依次作相邻对换，再把它与 $j$ 对换，然后把 $j$ 与前面
$2k$ 个元素依次作相邻对换，共作了 $2k+1$ 次相邻对换，每次都改变奇偶性，故
总效果是奇偶性改变。
\end{proof}

\begin{definition}\label{def:det-n}
设有 $n^2$ 个数排成 $n$ 行 $n$ 列，称
\[
  D=\begin{vmatrix}
      a_{11} & a_{12} & \cdots & a_{1n} \\
      a_{21} & a_{22} & \cdots & a_{2n} \\
      \vdots & \vdots & \ddots & \vdots \\
      a_{n1} & a_{n2} & \cdots & a_{nn}
    \end{vmatrix}
  =\sum_{p_1p_2\cdots p_n}(-1)^{t}\,a_{1p_1}a_{2p_2}\cdots a_{np_n}
\]
为\textbf{$n$ 阶行列式}，其中求和遍历列标排列 $p_1p_2\cdots p_n$ 的全部 $n!$
种取法，$t$ 为排列 $p_1p_2\cdots p_n$ 的逆序数。
\end{definition}

直观地说，$n$ 阶行列式是所有“取自不同行、不同列的 $n$ 个元素之积”的代数和：
每一项的行标按自然顺序排列，该项的符号由列标排列的奇偶性决定，偶排列取正号，
奇排列取负号。

按定义直接计算 $n!$ 项显然不现实，但某些特殊行列式可以立即得到结果。

\begin{example}\label{ex:upper-tri}
计算上三角行列式
$D=\begin{vmatrix} a_{11} & a_{12} & \cdots & a_{1n} \\ & a_{22} & \cdots & a_{2n} \\ & & \ddots & \vdots \\ 0 & & & a_{nn} \end{vmatrix}$
（主对角线以下元素均为 $0$）。
\end{example}
\begin{proof}[解]
考虑任一非零项 $a_{1p_1}a_{2p_2}\cdots a_{np_n}$。第 $n$ 行只有 $a_{nn}$ 可能
非零，故 $p_n=n$；第 $n-1$ 行只有 $a_{n-1,n-1}$ 与 $a_{n-1,n}$ 可能非零，
而列号 $n$ 已被占用，故 $p_{n-1}=n-1$；依此类推，唯一的非零项是
$a_{11}a_{22}\cdots a_{nn}$，其列标排列为自然排列，逆序数为 $0$。于是
\[
  D=a_{11}a_{22}\cdots a_{nn}.
\]
同理，下三角行列式与对角行列式也都等于主对角线元素之积。
\end{proof}

\begin{example}\label{ex:anti-diagonal}
计算 $n$ 阶行列式
\[
  D_n=\begin{vmatrix}
        0      &        &        & a_{1} \\
               &        & a_{2}  & \\
               & \iddots&        & \\
        a_{n}  &        &        & 0
      \end{vmatrix}
\]
（副对角线上元素依次为 $a_1,a_2,\cdots,a_n$，其余元素均为 $0$）。
\end{example}
\begin{proof}[解]
非零项必须取自副对角线，即第 $i$ 行取第 $n+1-i$ 列的元素，列标排列为
$n,\,n-1,\cdots,2,\,1$，其逆序数为
$t=0+1+2+\cdots+(n-1)=\dfrac{n(n-1)}{2}$。于是
\[
  D_n=(-1)^{\frac{n(n-1)}{2}}a_1a_2\cdots a_n .
\]
\end{proof}

\begin{insight}[本质：用三条要求刻画行列式]
$n!$ 项求和的定义不适合“理解”，它适合“证明”。行列式真正的本质是：它把
方阵对应成一个数，并且是满足下列三条自然要求的\textbf{唯一}函数——
（i）对每一行都是线性的（某行拆成两数之和，行列式随之拆开；某行提出公因子
$k$，行列式提出 $k$）；（ii）交换两行，行列式变号（特别地，两行相同则为
零）；（iii）单位矩阵的行列式为 $1$。先承认这三条，行列式的取值就被完全
确定，$n!$ 项的公式只是它的副产品。本章的所有性质、按行展开定理乃至克拉默
法则，都可以从这三条推出；在抽象代数中，这正是行列式的标准定义。
\end{insight}

\begin{qa}{定义中的每一项为什么必须取自“不同的行、不同的列”？符号为什么恰好由逆序数的奇偶决定？}
答：因为行列式“对每一行都是线性的”：把行列式按第 1 行、第 2 行……逐行
拆开，每一项中每一行恰好贡献一个因子，共 $n$ 个；再要求这 $n$ 个因子取自
不同列，才能既不重复又不遗漏地遍历所有取法，共 $n!$ 项。至于符号，它必须
满足“交换两行变号”：交换两行使每一项的列标排列恰好作一次对换，由定理
~\ref{thm:swap-parity}（一次对换改变奇偶性），所有项的符号一齐改变，整个
和恰好变号。换句话说，$(-1)^t$ 是让“变号”能够\textbf{整体一致}发生的
唯一可行的记账方式。
\end{qa}

\section{行列式的性质}\label{sec:det-properties}

把行列式 $D$ 的行与列互换（第 $i$ 行变为第 $i$ 列）所得的行列式称为 $D$ 的
\textbf{转置行列式}，记作 $D\trans$。下列性质对行与列同样成立，故只对行叙述。

\begin{property}[行列式与它的转置行列式相等]\label{prop:det-transpose}
$D=D\trans$。
\end{property}

按定义，$D\trans$ 的每一项可把因子按行标（即 $D$ 的列标）大小重新排序，
其符号由 $D$ 的行标排列决定；可以证明一个排列与其逆排列的逆序数奇偶性相同，
故 $D\trans$ 与 $D$ 的项一一对应且符号相同（证明从略）。由该性质，凡对行成立
的性质对列也成立。

\begin{property}[对换两行，行列式变号]\label{prop:det-swap}
对换行列式的两行（列），行列式变号。
\end{property}
\begin{proof}
对换第 $i,j$ 两行后，每一项的行标排列相当于原排列经过一次对换，由定理
\ref{thm:swap-parity}，每一项都改变符号，故行列式变号。
\end{proof}

\begin{property}[两行相同则为零]\label{prop:det-equal-rows}
若行列式有两行（列）完全相同，则此行列式等于 $0$。
\end{property}
\begin{proof}
对换相同的两行，行列式不变；但由性质~\ref{prop:det-swap} 它又应变号，即
$D=-D$，故 $D=0$。
\end{proof}

\begin{property}[行的公因子可提出]\label{prop:det-factor}
行列式某一行（列）的所有元素的公因子可以提到行列式记号外面，即
\[
  \begin{vmatrix}
    a_{11} & \cdots & a_{1n} \\
    \vdots &        & \vdots \\
    ka_{i1} & \cdots & ka_{in} \\
    \vdots &        & \vdots
  \end{vmatrix}
  =k
  \begin{vmatrix}
    a_{11} & \cdots & a_{1n} \\
    \vdots &        & \vdots \\
    a_{i1} & \cdots & a_{in} \\
    \vdots &        & \vdots
  \end{vmatrix}.
\]
\end{property}

由定义，每一项恰含第 $i$ 行的一个因子，提出公因子 $k$ 即得。

\begin{property}[两行成比例则为零]\label{prop:det-proportional}
若行列式有两行（列）的对应元素成比例，则此行列式等于 $0$。
\end{property}
\begin{proof}
设第 $j$ 行是第 $i$ 行的 $k$ 倍。由性质~\ref{prop:det-factor} 提出公因子 $k$
后，行列式有两行相同，再由性质~\ref{prop:det-equal-rows} 即得结论。
\end{proof}

\begin{property}[行可拆分]\label{prop:det-split}
若行列式的某一行的元素都是两数之和，则该行列式可拆成两个行列式之和，例如
\[
  \begin{vmatrix} a_{11} & a_{12} \\ b_1+c_1 & b_2+c_2 \end{vmatrix}
  =\begin{vmatrix} a_{11} & a_{12} \\ b_1 & b_2 \end{vmatrix}
  +\begin{vmatrix} a_{11} & a_{12} \\ c_1 & c_2 \end{vmatrix}.
\]
\end{property}
\begin{proof}
按定义，左边每一项 $a_{1p_1}(b_{p_2}+c_{p_2})$ 恰好拆成右边两个行列式对应
项之和，对应求和即得。
\end{proof}

\begin{property}[倍加不变]\label{prop:det-add}
把行列式某一行（列）的各元素乘以同一数 $k$ 后加到另一行（列）对应的元素上，
行列式不变。
\end{property}
\begin{proof}
以 $r_i+k r_j$ 代替 $r_i$（记号 $r_i$ 表示第 $i$ 行），由性质~\ref{prop:det-split}
与性质~\ref{prop:det-proportional}，
\[
  \begin{vmatrix} \vdots \\ r_i+kr_j \\ \vdots \end{vmatrix}
  =\begin{vmatrix} \vdots \\ r_i \\ \vdots \end{vmatrix}
  +\begin{vmatrix} \vdots \\ kr_j \\ \vdots \end{vmatrix}
  =D+0=D .\qedhere
\]
\end{proof}

\begin{qa}{为什么“某行加上另一行的倍数”不改变行列式，而“某行乘以 $k$”却要把 $k$ 提出行列式？}
答：代数上：倍加后按性质~\ref{prop:det-split} 拆开，第二项的两行成比例，
由性质~\ref{prop:det-proportional} 等于零；而数乘后\textbf{每一项}都恰含该
行的一个因子，$k$ 被所有项公共提取。几何上（图~\ref{fig:det-area}）：倍加
相当于把平行四边形沿一条边的方向“推斜”（剪切），底不变、高不变，面积当然
不变；数乘则真的把一条边拉长了 $k$ 倍，面积随之扩大 $k$ 倍。
\end{qa}

利用这些性质，可以把行列式化为上三角形式再计算，这是计算行列式的基本方法。

\begin{example}\label{ex:det-48}
计算四阶行列式
$D=\begin{vmatrix} 3 & 1 & 1 & 1 \\ 1 & 3 & 1 & 1 \\ 1 & 1 & 3 & 1 \\ 1 & 1 & 1 & 3 \end{vmatrix}$。
\end{example}
\begin{proof}[解]
\begin{align*}
  D &\xlongequal{r_i-r_1\ (i=2,3,4)}
  \begin{vmatrix} 3 & 1 & 1 & 1 \\ -2 & 2 & 0 & 0 \\ -2 & 0 & 2 & 0 \\ -2 & 0 & 0 & 2 \end{vmatrix}
  \xlongequal{c_1+c_2+c_3+c_4}
  \begin{vmatrix} 6 & 1 & 1 & 1 \\ 0 & 2 & 0 & 0 \\ 0 & 0 & 2 & 0 \\ 0 & 0 & 0 & 2 \end{vmatrix} \\
    &=6\times 2\times 2\times 2=48 ,
\end{align*}
其中 $c_j$ 表示第 $j$ 列，$r_i$ 表示第 $i$ 行。
\end{proof}

\section{行列式按行（列）展开}

高阶行列式可以通过“降阶”化为低阶行列式计算。

\begin{definition}\label{def:cofactor}
在 $n$ 阶行列式中，把元素 $a_{ij}$ 所在的第 $i$ 行和第 $j$ 列划去后，余下的
$n-1$ 阶行列式称为 $a_{ij}$ 的\textbf{余子式}，记作 $M_{ij}$；称
$A_{ij}=(-1)^{i+j}M_{ij}$ 为 $a_{ij}$ 的\textbf{代数余子式}。
\end{definition}

\begin{theorem}[展开定理]\label{thm:expansion}
行列式等于它的任一行（列）的各元素与其对应的代数余子式乘积之和，即
\[
  D=a_{i1}A_{i1}+a_{i2}A_{i2}+\cdots+a_{in}A_{in}\qquad (i=1,2,\cdots,n).
\]
\end{theorem}
\begin{proof}
先证特殊情形：第 $i$ 行除 $a_{ij}$ 外全为零。此时定义中的非零项必取 $a_{ij}$，
其余 $n-1$ 个因子取自划去第 $i$ 行、第 $j$ 列后的子式，该项的符号等于
$(-1)^{t}$，其中 $t$ 可分解为列标排列与行标排列两部分逆序之和，恰为
$(-1)^{i+j}M_{ij}$ 的符号，于是 $D=a_{ij}A_{ij}$。

一般情形：把第 $i$ 行写为 $n$ 个“只有一个非零元”的向量之和，由性质
\ref{prop:det-split} 拆成 $n$ 个上述特殊情形的行列式之和，即得
$D=\sum_{j=1}^{n}a_{ij}A_{ij}$。
\end{proof}

\begin{corollary}\label{cor:alien-cofactor}
行列式某一行（列）的元素与另一行（列）对应元素的代数余子式乘积之和为零，即
当 $i\neq k$ 时，
\[
  a_{i1}A_{k1}+a_{i2}A_{k2}+\cdots+a_{in}A_{kn}=0 .
\]
\end{corollary}
\begin{proof}
把 $D$ 的第 $k$ 行换成第 $i$ 行的元素，所得行列式 $D'$ 有两行相同，故 $D'=0$；
把 $D'$ 按第 $k$ 行展开，注意 $D'$ 的第 $k$ 行元素的代数余子式与 $D$ 的相同，
即得 $\sum_{j}a_{ij}A_{kj}=0$。
\end{proof}

\begin{example}\label{ex:expand-3}
按第一行展开计算
$D=\begin{vmatrix} 1 & 2 & 3 \\ 4 & 5 & 6 \\ 7 & 8 & 10 \end{vmatrix}$。
\end{example}
\begin{proof}[解]
\begin{align*}
  D &= 1\cdot\begin{vmatrix} 5 & 6 \\ 8 & 10 \end{vmatrix}
      -2\cdot\begin{vmatrix} 4 & 6 \\ 7 & 10 \end{vmatrix}
      +3\cdot\begin{vmatrix} 4 & 5 \\ 7 & 8 \end{vmatrix} \\
    &= 1\times(50-48)-2\times(40-42)+3\times(32-35)
     = 2+4-9=-3,
\end{align*}
与 \S 1.1 中对角线法则的结果一致。
\end{proof}

\begin{theorem}[范德蒙德行列式]\label{thm:vandermonde}
\[
  V_n=\begin{vmatrix}
    1      & 1      & \cdots & 1      \\
    x_1    & x_2    & \cdots & x_n    \\
    x_1^2  & x_2^2  & \cdots & x_n^2  \\
    \vdots & \vdots & \ddots & \vdots \\
    x_1^{n-1} & x_2^{n-1} & \cdots & x_n^{n-1}
  \end{vmatrix}
  =\prod_{1\le i<j\le n}(x_j-x_i).
\]
\end{theorem}
\begin{proof}
对 $n$ 作归纳。$n=2$ 时结论显然。设结论对 $n-1$ 成立。把 $V_n$ 看作最后一个
变量 $x_n$ 的多项式，其次数为 $n-1$。当 $x_n=x_i$（$i<n$）时行列式有两列
相同，故 $V_n=0$，即 $x_n-x_1,\cdots,x_n-x_{n-1}$ 都是该多项式的因式。这些
因式两两互素，其乘积的次数为 $n-1$，与 $V_n$ 关于 $x_n$ 的次数相同，故
\[
  V_n=c\,(x_n-x_1)(x_n-x_2)\cdots(x_n-x_{n-1}),
\]
其中 $c$ 是 $x_n^{n-1}$ 的系数。按最后一行展开，$x_n^{n-1}$ 只出现在代数余
子式为 $V_{n-1}(x_1,\cdots,x_{n-1})$ 的项中，故 $c=V_{n-1}$。由归纳假设
\[
  V_n=\prod_{i=1}^{n-1}(x_n-x_i)\cdot\prod_{1\le i<j\le n-1}(x_j-x_i)
     =\prod_{1\le i<j\le n}(x_j-x_i). \qedhere
\]
\end{proof}

例如
$\begin{vmatrix} 1 & 1 & 1 \\ 1 & 2 & 3 \\ 1 & 4 & 9 \end{vmatrix}
=(2-1)(3-1)(3-2)=2$。

\section{克拉默法则}\label{sec:cramer}

现在回到本章开始的线性方程组问题。

\begin{theorem}[克拉默法则]\label{thm:cramer}
含有 $n$ 个方程、$n$ 个未知数的线性方程组
\[
  \begin{cases}
    a_{11}x_1+a_{12}x_2+\cdots+a_{1n}x_n=b_1, \\
    a_{21}x_1+a_{22}x_2+\cdots+a_{2n}x_n=b_2, \\
    \quad\cdots\cdots                                         \\
    a_{n1}x_1+a_{n2}x_2+\cdots+a_{nn}x_n=b_n
  \end{cases}
\]
当其系数行列式
\[
  D=\begin{vmatrix} a_{11} & \cdots & a_{1n} \\ \vdots & & \vdots \\ a_{n1} & \cdots & a_{nn} \end{vmatrix}\neq 0
\]
时有唯一解，且解为
\[
  x_1=\frac{D_1}{D},\qquad x_2=\frac{D_2}{D},\qquad\cdots,\qquad x_n=\frac{D_n}{D},
\]
其中 $D_j$ 是把 $D$ 的第 $j$ 列换成常数项列 $b_1,b_2,\cdots,b_n$ 所得的行列式。
\end{theorem}
\begin{proof}
记方程组的系数矩阵为 $A$（矩阵的概念见第二章），$D_j=\sum_{k=1}^{n}b_kA_{kj}$
（把 $D_j$ 按第 $j$ 列展开）。

\textbf{解的存在性：}令 $x_j=D_j/D$。代入第 $i$ 个方程，利用推论~\ref{cor:alien-cofactor}
与展开定理，
\[
  \sum_{j=1}^{n}a_{ij}\frac{D_j}{D}
  =\frac{1}{D}\sum_{j=1}^{n}a_{ij}\sum_{k=1}^{n}b_kA_{kj}
  =\frac{1}{D}\sum_{k=1}^{n}b_k\sum_{j=1}^{n}a_{ij}A_{kj}
  =\frac{1}{D}\sum_{k=1}^{n}b_k\,\delta_{ik}D=b_i,
\]
其中 $\delta_{ik}=1$（$i=k$）或 $0$（$i\neq k$）。故所给 $x_j$ 确为解。

\textbf{解的唯一性：}设 $x_1,\cdots,x_n$ 为任一解。对每个 $j$，用代数余子式
$A_{mj}$（$m=1,\cdots,n$）乘第 $m$ 个方程后相加，$x_j$ 以外的未知数系数由
推论~\ref{cor:alien-cofactor} 知为零，得
\[
  \Big(\sum_{m=1}^{n}a_{mj}A_{mj}\Big)x_j
  =\sum_{m=1}^{n}b_mA_{mj},
  \quad\text{即}\quad
  D\,x_j=D_j,\quad x_j=\frac{D_j}{D}. \qedhere
\]
\end{proof}

\begin{corollary}\label{cor:homog}
若齐次线性方程组（$b_1=\cdots=b_n=0$）的系数行列式 $D\neq 0$，则它只有零解；
等价地，若齐次线性方程组有非零解，则 $D=0$。
\end{corollary}

\begin{example}\label{ex:cramer}
解线性方程组
\[
  \begin{cases}
    x_1+x_2+x_3=6, \\
    x_1+2x_2+3x_3=14, \\
    x_1+4x_2+9x_3=36.
  \end{cases}
\]
\end{example}
\begin{proof}[解]
系数行列式是范德蒙德行列式：
\[
  D=\begin{vmatrix} 1 & 1 & 1 \\ 1 & 2 & 3 \\ 1 & 4 & 9 \end{vmatrix}=2 .
\]
分别算得
\[
  D_1=\begin{vmatrix} 6 & 1 & 1 \\ 14 & 2 & 3 \\ 36 & 4 & 9 \end{vmatrix}=2,
  \qquad
  D_2=\begin{vmatrix} 1 & 6 & 1 \\ 1 & 14 & 3 \\ 1 & 36 & 9 \end{vmatrix}=4,
  \qquad
  D_3=\begin{vmatrix} 1 & 1 & 6 \\ 1 & 2 & 14 \\ 1 & 4 & 36 \end{vmatrix}=6,
\]
故 $x_1=D_1/D=1$，$x_2=D_2/D=2$，$x_3=D_3/D=3$。
\end{proof}

\begin{qa}{系数行列式 $D=0$ 时方程组会怎样？克拉默法则在实际计算中好用吗？}
答：克拉默法则只处理 $D\neq0$ 的情形，并断言此时解存在且唯一。当 $D=0$ 时
方程组可能无解、也可能有无穷多解，单靠行列式无法区分，这需要第四章的秩
判别法（比较 $r(A)$ 与 $r(A\mid b)$）。计算上，克拉默法则要算 $n+1$ 个 $n$
阶行列式，代价随 $n$ 急剧增长，实际数值求解都采用第四章的高斯消元法；克拉
默法则的价值在于给出解的\textbf{显式公式}与理论判据（例如立即得到推论
~\ref{cor:homog}，并在第二章给出可逆性判别）。
\end{qa}

\chapterexercises{一}

\begin{enumerate}[label=\arabic*., itemsep=2pt]
\item 计算下列行列式：
  \[
  \text{(1)}\ \begin{vmatrix} 2 & 0 & 1 \\ 1 & 3 & 5 \\ 4 & 1 & 2 \end{vmatrix};
  \qquad
  \text{(2)}\ \begin{vmatrix} a & b & c \\ b & c & a \\ c & a & b \end{vmatrix}.
  \]
\item 计算四阶行列式
  $\begin{vmatrix} 3 & 1 & 1 & 1 \\ 1 & 3 & 1 & 1 \\ 1 & 1 & 3 & 1 \\ 1 & 1 & 1 & 3 \end{vmatrix}$。
\item 计算四阶行列式
  $\begin{vmatrix} 1 & 2 & 0 & 0 \\ 3 & 4 & 0 & 0 \\ 0 & 0 & 5 & 6 \\ 0 & 0 & 7 & 8 \end{vmatrix}$。
\item 证明：若行列式的某一行是两组数的和，则行列式可按该行拆成两个行列式之和
  （性质~\ref{prop:det-split} 的一般形式）。
\item 利用性质证明性质~\ref{prop:det-proportional}（两行成比例的行列式为零）。
\item 用克拉默法则解方程组
  $\begin{cases} x_1+2x_2=5, \\ 3x_1+4x_2=9. \end{cases}$
\item 用克拉默法则解方程组
  $\begin{cases} x_1+x_2+x_3=3, \\ x_1+2x_2+3x_3=6, \\ 2x_1+x_2+x_3=4. \end{cases}$
\item $\lambda$ 取何值时，齐次线性方程组
  $\begin{cases} x_1+\lambda x_2=0, \\ \lambda x_1+x_2=0 \end{cases}$
  有非零解？
\item 证明：若 $n$ 元齐次线性方程组有非零解，则其系数行列式等于零。
\item 计算四阶范德蒙德行列式
  $\begin{vmatrix} 1 & 1 & 1 & 1 \\ 1 & 2 & 3 & 4 \\ 1 & 4 & 9 & 16 \\ 1 & 8 & 27 & 64 \end{vmatrix}$。
\end{enumerate}

\chapter{矩阵及其运算}

矩阵是线性代数的核心研究对象。本章介绍矩阵的概念与运算、逆矩阵、分块矩阵、
初等变换与初等矩阵，以及矩阵的秩。

\begin{insight}[观点：矩阵不只是数表，更是对空间的变换]
本书始终采用双重观点：矩阵既是一张按行列排好的数表，又是\textbf{对空间的一次
线性操作}——它把每个输入向量移到新的位置，且保持向量的加法与数乘（形象地说：
保持网格线平行、等距）。例如 $\begin{pmatrix} 0 & -1 \\ 1 & 0 \end{pmatrix}$ 把
平面绕原点逆时针旋转 $90^\circ$；第一章的行列式，正是这种操作对面积（体积）
的缩放因子。带着这个图像去读本章的每一条定义，“为什么这样规定”大多会
有着落；遇到代数上费解的结论，也先用变换的图像想一想。
\end{insight}

\section{矩阵的概念}

\begin{definition}\label{def:matrix}
由 $m\times n$ 个数 $a_{ij}$（$i=1,2,\cdots,m$；$j=1,2,\cdots,n$）排成的
$m$ 行 $n$ 列的数表
\[
  A=\begin{pmatrix}
    a_{11} & a_{12} & \cdots & a_{1n} \\
    a_{21} & a_{22} & \cdots & a_{2n} \\
    \vdots & \vdots & \ddots & \vdots \\
    a_{m1} & a_{m2} & \cdots & a_{mn}
  \end{pmatrix}
\]
称为 $m$ 行 $n$ 列矩阵，简称 $m\times n$ \textbf{矩阵}。这 $mn$ 个数称为矩阵
$A$ 的元素，$a_{ij}$ 位于第 $i$ 行第 $j$ 列。元素为实数的矩阵称为实矩阵，
本书中的矩阵除特别说明外均指实矩阵。
\end{definition}

矩阵常用大写字母 $A,B,C,\cdots$ 表示，也可记作 $A=(a_{ij})_{m\times n}$ 或
$A_{m\times n}$。行数与列数相等的矩阵称为\textbf{方阵}，$n$ 行 $n$ 列的方阵
简称为 $n$ 阶方阵。

两个矩阵若行数、列数分别相同，则称它们为\textbf{同型矩阵}；若两个同型矩阵的
对应元素都相等，则称两矩阵\textbf{相等}。

常见的特殊矩阵有：
\begin{itemize}[itemsep=1pt]
  \item \textbf{零矩阵} $O$：所有元素都是 $0$ 的矩阵；
  \item \textbf{单位矩阵} $E$（或记作 $E_n$）：主对角线上元素都是 $1$、其余元素
        都是 $0$ 的 $n$ 阶方阵；
  \item \textbf{对角矩阵}（数量矩阵）：主对角线以外的元素都是 $0$ 的方阵，
        记作 $\operatorname{diag}(d_1,d_2,\cdots,d_n)$；当 $d_1=\cdots=d_n=k$ 时
        称为数量矩阵 $kE$；
  \item \textbf{上（下）三角矩阵}：主对角线以下（上）的元素都是 $0$ 的方阵；
  \item \textbf{对称矩阵}：满足 $A\trans=A$ 的方阵；\textbf{反对称矩阵}：
        满足 $A\trans=-A$ 的方阵。
\end{itemize}

只有一行（列）的矩阵分别称为\textbf{行向量}与\textbf{列向量}。

必须强调：方阵与行列式是两个不同的概念。$n$ 阶方阵是一个数表，而 $n$ 阶
行列式是按第一章规则由该数表算出的一个数，记作 $|A|$ 或 $\det A$。

\section{矩阵的运算}\label{sec:matrix-ops}

\textbf{（1）加法}：设 $A=(a_{ij})$、$B=(b_{ij})$ 是两个 $m\times n$ 矩阵，
定义
\[
  A+B=(a_{ij}+b_{ij})_{m\times n}.
\]
只有同型矩阵才能相加。加法满足交换律与结合律；$A+O=A$；记 $-A=(-a_{ij})$，
则 $A+(-A)=O$，矩阵的减法定义为 $A-B=A+(-B)$。

\textbf{（2）数乘}：设 $\lambda$ 为数，定义 $\lambda A=(\lambda a_{ij})$。
数乘满足结合律以及对矩阵加法、数加法的分配律。

\textbf{（3）乘法}：设 $A=(a_{ij})$ 是 $m\times s$ 矩阵，$B=(b_{ij})$ 是
$s\times n$ 矩阵，定义 $AB$ 为 $m\times n$ 矩阵 $(c_{ij})$，其中
\begin{equation}\label{eq:matrix-multiply}
  c_{ij}=a_{i1}b_{1j}+a_{i2}b_{2j}+\cdots+a_{is}b_{sj}
        =\sum_{k=1}^{s}a_{ik}b_{kj}.
\end{equation}
即乘积的第 $i$ 行第 $j$ 列元素等于左矩阵第 $i$ 行与右矩阵第 $j$ 列对应元素
乘积之和。只有左矩阵的列数等于右矩阵的行数时，两矩阵才能相乘。

\begin{insight}[乘法的几何意义：列的组合与变换的复合]
把 $A$ 的第 $j$ 列记作 $a_j$，则规则 \eqref{eq:matrix-multiply} 等价于
\[
  Ax=x_1a_1+x_2a_2+\cdots+x_sa_s ,
\]
即：\textbf{$Ax$ 就是把 $x$ 的分量当作权重，对 $A$ 的各列作线性组合}。于是
一个矩阵完整描述了一个线性变换：把标准基 $\varepsilon_j$ 输入，输出恰是 $A$
的第 $j$ 列——\textbf{矩阵的列，就是基向量被变换后到达的位置}。矩阵乘积
$BA$ 则是“先做 $A$，再做 $B$”的复合变换：这解释了结合律为何自然成立，
也预示了交换律为何一般不成立（先旋转后剪切，与先剪切后旋转，结果不同）。
\end{insight}

\begin{figure}[htbp]
\centering
\begin{tikzpicture}[>=Stealth, scale=0.85]
  \begin{scope}
    \draw[gray!50] (-0.5,-0.5) grid[step=0.5] (2.5,2.5);
    \draw[->] (-0.6,0) -- (2.8,0);
    \draw[->] (0,-0.6) -- (0,2.8);
    \draw[very thick, blue!60!black, ->] (0,0) -- (1,0) node[below right] {$\varepsilon_1$};
    \draw[very thick, red!70!black, ->] (0,0) -- (0,1) node[left] {$\varepsilon_2$};
    \node[below] at (1,-1.15) {变换前：均匀方格};
  \end{scope}
  \draw[->, very thick, gray] (3.6,1) -- (5.2,1)
        node[midway, above] {$A=\begin{pmatrix} 1 & 1 \\ 0 & 1 \end{pmatrix}$};
  \begin{scope}[xshift=6.4cm]
    \foreach \y in {0.5,1,1.5,2}
      \draw[gray!50] (-0.3,\y) -- (3.4,\y);
    \foreach \k in {0,0.5,1,1.5,2,2.5}
      \draw[gray!50] (\k,0) -- (\k,0) ++(2.6,2.6);
    \draw[->] (-0.6,0) -- (3.9,0);
    \draw[->] (0,-0.6) -- (0,2.8);
    \draw[very thick, blue!60!black, ->] (0,0) -- (1,0) node[below] {$(1,0)\trans$};
    \draw[very thick, red!70!black, ->] (0,0) -- (1,1) node[above right] {$(1,1)\trans$};
    \node[below] at (1.6,-1.15) {变换后：方格被“剪切”推斜};
  \end{scope}
\end{tikzpicture}
\caption{剪切变换 $A=\begin{pmatrix} 1 & 1 \\ 0 & 1 \end{pmatrix}$：
$\varepsilon_1$ 不动，$\varepsilon_2$ 移到 $(1,1)\trans$，整张方格被推斜，
但网格线仍平行、等距。面积因子 $|A|=1$，与例~\ref{ex:power} 中 $A$ 的幂的
行为一致。}
\label{fig:shear}
\end{figure}

矩阵乘法满足结合律 $(AB)C=A(BC)$、左分配律 $A(B+C)=AB+AC$、右分配律
$(B+C)A=BA+CA$，以及 $\lambda(AB)=(\lambda A)B=A(\lambda B)$；对任意
$m\times n$ 矩阵 $A$，有 $E_mA=A=AE_n$。

\begin{remark}\label{rmk:no-commute}
与数的乘法不同，矩阵乘法\textbf{不满足交换律}。首先，当 $A$ 的列数不等于 $B$
的行数时 $AB$ 有意义而 $BA$ 未必有意义；其次，即使 $AB$ 与 $BA$ 都有意义，
它们也未必是同型矩阵；最后，即使同型，也未必相等。
\end{remark}

\begin{example}
设 $A=\begin{pmatrix} 1 & 0 \\ 1 & 1 \end{pmatrix}$，
$B=\begin{pmatrix} 1 & 1 \\ 0 & 1 \end{pmatrix}$，求 $AB$ 与 $BA$。
\end{example}
\begin{proof}[解]
\[
  AB=\begin{pmatrix} 1 & 1 \\ 1 & 2 \end{pmatrix},
  \qquad
  BA=\begin{pmatrix} 2 & 1 \\ 1 & 1 \end{pmatrix},
\]
可见 $AB\neq BA$。
\end{proof}

由交换律不成立还可推出：一般的“平方差公式” $(A+B)(A-B)=A^2-B^2$ 对矩阵
不成立（右边会多出交叉项），只有当 $AB=BA$ 时才有 $(A+B)(A-B)=A^2-B^2$。

\begin{remark}\label{rmk:no-cancel}
由 $AB=O$ 一般不能推出 $A=O$ 或 $B=O$，从而矩阵乘法\textbf{不满足消去律}：
由 $AB=AC$ 且 $A\neq O$ 一般不能推出 $B=C$。
\end{remark}

\begin{example}
设 $A=\begin{pmatrix} 1 & 1 \\ 1 & 1 \end{pmatrix}$，
$B=\begin{pmatrix} 1 & -1 \\ -1 & 1 \end{pmatrix}$，则
\[
  AB=\begin{pmatrix} 1-1 & -1+1 \\ 1-1 & -1+1 \end{pmatrix}
    =\begin{pmatrix} 0 & 0 \\ 0 & 0 \end{pmatrix},
\]
而 $A,B$ 都不是零矩阵。
\end{example}

\begin{qa}{$AB=O$ 却 $A\neq O$、$B\neq O$，几何上是怎么回事？消去律为什么失效？}
答：把 $B$ 按列看：$AB=(Ab_1,\cdots,Ab_n)$，故 $AB=O$ 说明\textbf{$B$ 的每一列
都被 $A$ 变换成了零向量}。只要 $A$ 把某些方向“压扁”（例如
$\operatorname{diag}(1,0)$ 把整个 $x_2$ 轴压到原点），取 $B$ 的列全部落在被
压扁的方向上，乘积就是零，而两个因子都非零。消去律失效是同一件事：
$AB=AC$ 即 $A(B-C)=O$，$B-C$ 的列可能恰好都落在被压扁的方向上，未必为零。
\end{qa}

\textbf{（4）转置}：把矩阵 $A$ 的行换成同序号的列所得的矩阵称为 $A$ 的
\textbf{转置矩阵}，记作 $A\trans$。转置满足
\[
  (A\trans)\trans=A,\qquad (A+B)\trans=A\trans+B\trans,\qquad
  (\lambda A)\trans=\lambda A\trans,\qquad (AB)\trans=B\trans A\trans .
\]
最后一个等式证明如下：$(AB)\trans$ 的 $(i,j)$ 元为 $(AB)_{ji}
=\sum_k a_{jk}b_{ki}$，而 $B\trans A\trans$ 的 $(i,j)$ 元为
$\sum_k (B\trans)_{ik}(A\trans)_{kj}=\sum_k b_{ki}a_{jk}$，两者相同。

\textbf{（5）方阵的幂与多项式}：设 $A$ 为 $n$ 阶方阵，定义
$A^k=A\cdots A$（$k$ 个相乘，$k$ 为正整数），并约定 $A^0=E$。由于乘法结合律
成立，有 $A^kA^l=A^{k+l}$，$(A^k)^l=A^{kl}$；但由于交换律不成立，
$(A+B)^2=A^2+AB+BA+B^2$ 而非 $A^2+2AB+B^2$。

\begin{example}\label{ex:power}
设 $A=\begin{pmatrix} 1 & 1 \\ 0 & 1 \end{pmatrix}$，证明
$A^n=\begin{pmatrix} 1 & n \\ 0 & 1 \end{pmatrix}$（$n$ 为正整数）。
\end{example}
\begin{proof}
$n=1$ 时显然。设结论对 $n$ 成立，则
\[
  A^{n+1}=A^nA
  =\begin{pmatrix} 1 & n \\ 0 & 1 \end{pmatrix}
   \begin{pmatrix} 1 & 1 \\ 0 & 1 \end{pmatrix}
  =\begin{pmatrix} 1 & n+1 \\ 0 & 1 \end{pmatrix},
\]
由归纳法结论成立。
\end{proof}

设 $f(x)=a_mx^m+\cdots+a_1x+a_0$ 为多项式，$A$ 为方阵，称
$f(A)=a_mA^m+\cdots+a_1A+a_0E$ 为 $A$ 的\textbf{矩阵多项式}。

\begin{example}
设 $A$ 为对称矩阵，$B$ 为反对称矩阵，$AB$ 何时为对称矩阵？
\end{example}
\begin{proof}[解]
$(AB)\trans=B\trans A\trans=-BA$。要使 $AB$ 对称，即 $-BA=AB$，等价于
$AB=-BA=BA$，即 $A$ 与 $B$ 可交换。
\end{proof}

\section{逆矩阵}\label{sec:inverse}

数的乘法有逆运算除法，矩阵乘法的“逆运算”通过逆矩阵实现。

\begin{definition}\label{def:inverse}
设 $A$ 为 $n$ 阶方阵。若存在 $n$ 阶方阵 $B$，使
\[
  AB=BA=E,
\]
则称 $A$ \textbf{可逆}，并称 $B$ 为 $A$ 的\textbf{逆矩阵}。
\end{definition}

首先，若 $A$ 可逆，则逆矩阵唯一：设 $B,C$ 都是 $A$ 的逆矩阵，则
$B=BE=B(AC)=(BA)C=EC=C$。因此把 $A$ 的逆矩阵记作 $A^{-1}$。

方阵的逆满足（设下列方阵都可逆）：
\[
  (A^{-1})^{-1}=A,\qquad (AB)^{-1}=B^{-1}A^{-1},\qquad
  (\lambda A)^{-1}=\lambda^{-1}A^{-1},\qquad
  (A\trans)^{-1}=(A^{-1})\trans,
\]
且 $|A^{-1}|=|A|^{-1}$。其中 $(AB)^{-1}=B^{-1}A^{-1}$ 的验证：
$(B^{-1}A^{-1})(AB)=B^{-1}(A^{-1}A)B=B^{-1}B=E$，同理可验另一侧。

下面给出可逆的充要条件与逆矩阵的计算公式。对 $n$ 阶方阵 $A$，把 $|A|$ 的
各个元素的代数余子式按如下方式排成的 $n$ 阶方阵
\[
  A^*=\begin{pmatrix}
    A_{11} & A_{21} & \cdots & A_{n1} \\
    A_{12} & A_{22} & \cdots & A_{n2} \\
    \vdots & \vdots & \ddots & \vdots \\
    A_{1n} & A_{2n} & \cdots & A_{nn}
  \end{pmatrix}
\]
称为 $A$ 的\textbf{伴随矩阵}（注意代数余子式的行、列位置恰好互换）。

\begin{lemma}\label{lem:adjugate}
设 $A$ 为 $n$ 阶方阵，$A^*$ 为其伴随矩阵，则
\[
  AA^*=A^*A=|A|E .
\]
\end{lemma}
\begin{proof}
$AA^*$ 的 $(i,k)$ 元为 $\sum_{j=1}^{n}a_{ij}A_{kj}$。当 $k=i$ 时由第一章
展开定理（定理~\ref{thm:expansion}）该和等于 $|A|$；当 $k\neq i$ 时由其推论
（推论~\ref{cor:alien-cofactor}）该和等于 $0$。故 $AA^*=|A|E$。同理由按列
展开的对应结论得 $A^*A=|A|E$。
\end{proof}

\begin{theorem}\label{thm:inverse-det}
方阵 $A$ 可逆的充分必要条件是 $|A|\neq 0$，且当 $|A|\neq 0$ 时
\[
  A^{-1}=\frac{1}{|A|}A^* .
\]
\end{theorem}
\begin{proof}
若 $|A|\neq 0$，由引理~\ref{lem:adjugate} 得
$A\big(\tfrac{1}{|A|}A^*\big)=\big(\tfrac{1}{|A|}A^*\big)A=E$，故 $A$ 可逆且
$A^{-1}=\tfrac{1}{|A|}A^*$。反之，若 $A$ 可逆，设 $AB=E$，两边取行列式得
$|A||B|=1$，故 $|A|\neq 0$。
\end{proof}

$|A|\neq 0$ 的方阵称为\textbf{非奇异矩阵}，$|A|=0$ 的方阵称为\textbf{奇异矩阵}。
上述定理表明：可逆 $\iff$ 非奇异。

\begin{insight}[可逆的几何意义：变换能够“原路退回”]
可逆矩阵是一个\textbf{不丢失信息}的变换：不同的输入被送到不同的输出，因而
存在逆变换把一切还原。行列式恰好是判据：$|A|\neq0$（面积因子非零）时空间
没有被压扁，变换是一一对应；$|A|=0$ 时整个空间被压缩到更低维，无数不同的
输入共享同一个输出，还原不再可能。
\end{insight}

\begin{figure}[htbp]
\centering
\begin{tikzpicture}[>=Stealth, scale=0.9]
  \begin{scope}
    \draw[gray!50] (-0.5,-0.5) grid[step=0.5] (2.5,2.5);
    \draw[->] (-0.6,0) -- (2.8,0);
    \draw[->] (0,-0.6) -- (0,2.8);
    \fill[blue!60!black] (0.5,1.5) circle (2.2pt);
    \fill[blue!60!black] (1.5,1.0) circle (2.2pt);
    \fill[blue!60!black] (2.0,0.5) circle (2.2pt);
    \node[below] at (1,-1.05) {三个不同的输入向量};
  \end{scope}
  \draw[->, very thick, gray] (3.6,1) -- (5.4,1)
        node[midway, above] {$A=\operatorname{diag}(1,0)$};
  \begin{scope}[xshift=6.6cm]
    \draw[->] (-0.6,0) -- (3.2,0);
    \draw[very thick, red!70] (-0.3,0) -- (3.0,0);
    \fill[red!70!black] (0.5,0) circle (2.2pt);
    \fill[red!70!black] (1.5,0) circle (2.2pt);
    \fill[red!70!black] (2.0,0) circle (2.2pt);
    \node[below] at (1.3,-0.85) {三个输出全部压在 $y_1$ 轴上};
  \end{scope}
\end{tikzpicture}
\caption{降维压缩不可逆：$|A|=0$ 时，无数不同的输入被送到同一点，无法区分
“谁是谁”，因此不存在逆变换。可逆 $\iff$ 变换不把空间压扁。}
\label{fig:collapse}
\end{figure}

\begin{qa}{$|A|=0$、“方程组解不唯一”、“矩阵不可逆”、“秩亏”、“列向量线性相关”之间是什么关系？}
答：对 $n$ 阶方阵，这些说法互相等价：$|A|=0$ $\iff$ $A$ 不可逆 $\iff$
$\rank(A)<n$ $\iff$ 齐次方程组 $Ax=0$ 有非零解 $\iff$ $A$ 的列向量线性相关
（第三章）$\iff$ 变换把 $n$ 维空间压扁。读者目前可以完全看懂前三个等价；
后两个分别在第三章、第四章展开。把这张“等价网”记在心里，线性代数的各章
就串成了一条线。
\end{qa}

\begin{corollary}\label{cor:one-sided}
若 $AB=E$（或 $BA=E$），则 $B=A^{-1}$。
\end{corollary}
\begin{proof}
$|A||B|=|E|=1$，故 $|A|\neq 0$，$A$ 可逆，于是
$B=A^{-1}(AB)=A^{-1}E=A^{-1}$。
\end{proof}

\begin{example}\label{ex:inv-2}
设 $A=\begin{pmatrix} a & b \\ c & d \end{pmatrix}$，$|A|=ad-bc\neq 0$，求
$A^{-1}$。
\end{example}
\begin{proof}[解]
四个代数余子式为 $A_{11}=d$，$A_{12}=-c$，$A_{21}=-b$，$A_{22}=a$，故
\[
  A^{-1}=\frac{1}{ad-bc}\begin{pmatrix} d & -b \\ -c & a \end{pmatrix}.
\]
\end{proof}

\begin{example}\label{ex:inv-3}
求 $A=\begin{pmatrix} 1 & 2 & 3 \\ 2 & 2 & 3 \\ 1 & 1 & 1 \end{pmatrix}$ 的
逆矩阵。
\end{example}
\begin{proof}[解]
先算各代数余子式：
\[
  \begin{aligned}
    & A_{11}=-1,\quad && A_{12}=1,\quad && A_{13}=0, \\
    & A_{21}=1,\quad  && A_{22}=-2,\quad && A_{23}=1, \\
    & A_{31}=0,\quad  && A_{32}=3,\quad  && A_{33}=-2,
  \end{aligned}
\]
按第一行展开得 $|A|=1\times(-1)+2\times 1+3\times 0=1\neq 0$，故
\[
  A^*=\begin{pmatrix} -1 & 1 & 0 \\ 1 & -2 & 3 \\ 0 & 1 & -2 \end{pmatrix},
  \qquad
  A^{-1}=A^*=\begin{pmatrix} -1 & 1 & 0 \\ 1 & -2 & 3 \\ 0 & 1 & -2 \end{pmatrix}.
\]
可直接验证 $AA^{-1}=E$。
\end{proof}

\begin{example}\label{ex:mat-eq}
设 $A=\begin{pmatrix} 1 & 2 \\ 1 & 3 \end{pmatrix}$，
$B=\begin{pmatrix} 1 & 1 \\ 2 & 3 \end{pmatrix}$，解矩阵方程 $AX=B$。
\end{example}
\begin{proof}[解]
$|A|=1\neq 0$，故 $X=A^{-1}B$。由例~\ref{ex:inv-2}，
$A^{-1}=\begin{pmatrix} 3 & -2 \\ -1 & 1 \end{pmatrix}$，于是
\[
  X=\begin{pmatrix} 3 & -2 \\ -1 & 1 \end{pmatrix}
    \begin{pmatrix} 1 & 1 \\ 2 & 3 \end{pmatrix}
   =\begin{pmatrix} -1 & -3 \\ 1 & 2 \end{pmatrix}.
\]
\end{proof}

\section{分块矩阵}\label{sec:block}

对于行数、列数较高的矩阵，常用若干条贯穿始终的横线与纵线把它分成许多小块，
每一块是一个低阶矩阵，把每一小块当作一个元素处理，称为\textbf{矩阵的分块}。
分块矩阵的加法、数乘按块进行；分块矩阵相乘时，把小块当作元素按规则
\eqref{eq:matrix-multiply} 计算，但要求左矩阵的列分法与右矩阵的行分法一致。

\begin{example}
设
\[
  A=\begin{pmatrix} 1 & 0 & 0 \\ 0 & 1 & 0 \\ 0 & 0 & 0 \end{pmatrix}
\]
按 $\begin{pmatrix} E_2 & O \\ O & 0 \end{pmatrix}$ 分块，说明把 $A$ 左乘到
$3\times n$ 矩阵 $B$ 上的效果。
\end{example}
\begin{proof}[解]
把 $B$ 按行分块为 $B=\begin{pmatrix} B_1 \\ B_2 \\ B_3 \end{pmatrix}$，则
\[
  AB=\begin{pmatrix} E_2 & O \\ O & 0 \end{pmatrix}
      \begin{pmatrix} B_1 \\ B_2 \\ B_3 \end{pmatrix}
    =\begin{pmatrix} B_1 \\ B_2 \\ O \end{pmatrix},
\]
即 $A$ 把 $B$ 的第三行“清零”，前两行保持不变。
\end{proof}

分块对角矩阵（准对角矩阵）
$\operatorname{diag}(A_1,\cdots,A_s)$（其中 $A_i$ 都是方阵）的运算特别简单：
\[
  \operatorname{diag}(A_1,\cdots,A_s)\cdot\operatorname{diag}(B_1,\cdots,B_s)
  =\operatorname{diag}(A_1B_1,\cdots,A_sB_s),
\]
当每个 $A_i$ 都可逆时，
\[
  \operatorname{diag}(A_1,\cdots,A_s)^{-1}
  =\operatorname{diag}(A_1^{-1},\cdots,A_s^{-1}).
\]

\begin{example}
利用分块求
$A=\begin{pmatrix} 1 & 2 & 0 & 0 \\ 2 & 5 & 0 & 0 \\ 0 & 0 & 2 & 1 \\ 0 & 0 & 1 & 1 \end{pmatrix}$
的逆矩阵。
\end{example}
\begin{proof}[解]
$A=\operatorname{diag}(A_1,A_2)$，其中
$A_1=\begin{pmatrix} 1 & 2 \\ 2 & 5 \end{pmatrix}$，
$A_2=\begin{pmatrix} 2 & 1 \\ 1 & 1 \end{pmatrix}$，$|A_1|=|A_2|=1$。由例
~\ref{ex:inv-2}，
\[
  A_1^{-1}=\begin{pmatrix} 5 & -2 \\ -2 & 1 \end{pmatrix},
  \qquad
  A_2^{-1}=\begin{pmatrix} 1 & -1 \\ -1 & 2 \end{pmatrix},
\]
故
\[
  A^{-1}=\begin{pmatrix}
    5 & -2 & 0 & 0 \\ -2 & 1 & 0 & 0 \\ 0 & 0 & 1 & -1 \\ 0 & 0 & -1 & 2
  \end{pmatrix}.
\]
\end{proof}

分块矩阵转置时，除块本身要转置外，块的位置也要转置：若
$A=\begin{pmatrix} A_{11} & A_{12} \\ A_{21} & A_{22} \end{pmatrix}$，则
$A\trans=\begin{pmatrix} A_{11}\trans & A_{21}\trans \\ A_{12}\trans & A_{22}\trans \end{pmatrix}$。

\section{矩阵的初等变换与初等矩阵}\label{sec:elementary}

\begin{definition}\label{def:elementary-op}
下列三种变换称为矩阵的\textbf{初等行变换}：
\begin{enumerate}[label=(\arabic*), itemsep=0pt]
  \item 对换两行（对换 $i,j$ 两行，记作 $r_i\leftrightarrow r_j$）；
  \item 以非零数 $k$ 乘某一行（记作 $r_i\times k$）；
  \item 把某一行乘以数 $k$ 后加到另一行上（记作 $r_i+kr_j$）。
\end{enumerate}
把上述定义中的“行”换成“列”，即得初等列变换（记号相应地把 $r$ 换成 $c$）。
初等行变换与初等列变换统称\textbf{初等变换}。
\end{definition}

若矩阵 $A$ 经有限次初等变换变为 $B$，则称 $A$ 与 $B$ \textbf{等价}，记作
$A\sim B$。矩阵间的等价关系满足反身性、对称性与传递性。

\begin{definition}\label{def:echelon}
称满足下列条件的矩阵为\textbf{行阶梯形矩阵}：
（i）零行（元素全为零的行）位于非零行的下方；
（ii）各非零行的第一个非零元（称为主元）的列标随着行标的增大而严格增大。
若行阶梯形矩阵的主元都是 $1$，且主元所在列的其余元素都是 $0$，则称为
\textbf{行最简形矩阵}。
\end{definition}

\begin{theorem}\label{thm:echelon}
任何矩阵都可经有限次初等行变换化为行阶梯形，进而化为行最简形。
\end{theorem}
\begin{proof}
设 $A$ 的第 1 列有非零元。若 $a_{11}=0$ 而 $a_{i1}\neq 0$，先作
$r_1\leftrightarrow r_i$，使 $a_{11}\neq 0$；再对 $i=2,\cdots,m$ 作
$r_i-\frac{a_{i1}}{a_{11}}r_1$，把第一列主元下方化为 $0$。对右下角的
$(m-1)\times(n-1)$ 子块重复上述步骤，直到剩下的子块全是零行或没有剩余列。
由于每一步非零行数或已处理列数至少增加一个，过程有限步终止，得行阶梯形。
再用 $\frac{1}{a}\times r$ 把各主元化为 $1$，并自下而上用
$r_i-a_{ij}r_j$ 把主元上方化为 $0$，即得行最简形。
\end{proof}

\begin{definition}\label{def:elementary-matrix}
由单位矩阵 $E$ 经一次初等变换得到的矩阵称为\textbf{初等矩阵}，共三类：
$E(i,j)$（对换 $i,j$ 两行）、$E(i(k))$（第 $i$ 行乘 $k\neq 0$）、
$E(ij(k))$（第 $j$ 行乘 $k$ 加到第 $i$ 行上）。
\end{definition}

\begin{theorem}\label{thm:elem-mult}
设 $A$ 为 $m\times n$ 矩阵。对 $A$ 施行一次初等\textbf{行}变换，相当于在
$A$ 的\textbf{左边}乘以相应的 $m$ 阶初等矩阵；对 $A$ 施行一次初等\textbf{列}
变换，相当于在 $A$ 的\textbf{右边}乘以相应的 $n$ 阶初等矩阵。
\end{theorem}

以 $E(12)$ 左乘 $\begin{pmatrix} a_{11} & a_{12} \\ a_{21} & a_{22} \end{pmatrix}$
得 $\begin{pmatrix} a_{21} & a_{22} \\ a_{11} & a_{12} \end{pmatrix}$，恰好是
对换一、二两行的结果，其余情形类似（验证留给读者）。

初等矩阵都是可逆的，且其逆矩阵仍是同一类型的初等矩阵：
\[
  E(i,j)^{-1}=E(i,j),\qquad
  E(i(k))^{-1}=E\big(i(\tfrac{1}{k})\big),\qquad
  E(ij(k))^{-1}=E(ij(-k)).
\]

\begin{theorem}\label{thm:elem-product}
设 $A$ 为可逆方阵，则 $A$ 可以表示为有限个初等矩阵的乘积。
\end{theorem}
\begin{proof}
$A$ 可逆故 $|A|\neq 0$。由定理~\ref{thm:echelon}，$A$ 可经有限次初等行变换
化为行最简形 $R$，且每一步对应左乘一个初等矩阵，即存在初等矩阵
$P_1,\cdots,P_t$ 使 $P_t\cdots P_1A=R$。$|R|=|P_t\cdots P_1|\,|A|\neq 0$，
故 $R$ 没有零行，行最简形且满秩的方阵只能是 $E$，于是
$P_t\cdots P_1A=E$，即
\[
  A=P_1^{-1}\cdots P_t^{-1}E=P_1^{-1}\cdots P_t^{-1},
\]
是有限个初等矩阵之积。
\end{proof}

由此得到用初等变换求逆矩阵的方法。设 $A$ 可逆，由上述定理存在初等矩阵
$P$（若干个初等矩阵之积）使 $PA=E$，即 $P=A^{-1}$。对分块矩阵
$(A\mid E)$ 施行同样的初等行变换：
\[
  P(A\mid E)=(PA\mid PE)=(E\mid P)=(E\mid A^{-1}),
\]
即：\textbf{对 $(A\mid E)$ 施行初等行变换，当左半部分化为 $E$ 时，右半部分
就是 $A^{-1}$}；若左半部分出现零行，则 $A$ 不可逆。

\begin{example}\label{ex:inv-elementary}
用初等行变换求
$A=\begin{pmatrix} 2 & 1 & 1 \\ 1 & 2 & 1 \\ 1 & 1 & 2 \end{pmatrix}$ 的逆。
\end{example}
\begin{proof}[解]
\begin{align*}
  (A\mid E)
  &=\begin{pmatrix} 2 & 1 & 1 &\vline& 1 & 0 & 0 \\
                    1 & 2 & 1 &\vline& 0 & 1 & 0 \\
                    1 & 1 & 2 &\vline& 0 & 0 & 1 \end{pmatrix}
   \xrightarrow[\ r_3-2r_1\ ]{r_1\leftrightarrow r_3,\ r_2-r_1}
    \begin{pmatrix} 1 & 1 & 2 &\vline& 0 & 0 & 1 \\
                    0 & 1 & -1 &\vline& 0 & 1 & -1 \\
                    0 & -1 & -3 &\vline& 1 & 0 & -2 \end{pmatrix} \\
  &\xrightarrow{\ r_3+r_2\ }
    \begin{pmatrix} 1 & 1 & 2 &\vline& 0 & 0 & 1 \\
                    0 & 1 & -1 &\vline& 0 & 1 & -1 \\
                    0 & 0 & -4 &\vline& 1 & 1 & -3 \end{pmatrix}
   \xrightarrow{\ r_3\times(-\tfrac14)\ }
    \begin{pmatrix} 1 & 1 & 2 &\vline& 0 & 0 & 1 \\
                    0 & 1 & -1 &\vline& 0 & 1 & -1 \\
                    0 & 0 & 1 &\vline& -\tfrac14 & -\tfrac14 & \tfrac34 \end{pmatrix} \\
  &\xrightarrow[\ r_1-2r_3\ ]{r_2+r_3}
    \begin{pmatrix} 1 & 1 & 0 &\vline& \tfrac12 & \tfrac12 & -\tfrac12 \\
                    0 & 1 & 0 &\vline& -\tfrac14 & \tfrac34 & -\tfrac14 \\
                    0 & 0 & 1 &\vline& -\tfrac14 & -\tfrac14 & \tfrac34 \end{pmatrix}
   \xrightarrow{\ r_1-r_2\ }
    \begin{pmatrix} 1 & 0 & 0 &\vline& \tfrac34 & -\tfrac14 & -\tfrac14 \\
                    0 & 1 & 0 &\vline& -\tfrac14 & \tfrac34 & -\tfrac14 \\
                    0 & 0 & 1 &\vline& -\tfrac14 & -\tfrac14 & \tfrac34 \end{pmatrix},
\end{align*}
故
\[
  A^{-1}=\frac14\begin{pmatrix} 3 & -1 & -1 \\ -1 & 3 & -1 \\ -1 & -1 & 3 \end{pmatrix}.
\]
\end{proof}

\section{矩阵的秩}\label{sec:rank}

\begin{definition}\label{def:rank}
在 $m\times n$ 矩阵 $A$ 中任取 $k$ 行 $k$ 列（$k\le\min(m,n)$），位于这些行、
列交叉处的 $k^2$ 个元素按原次序组成的 $k$ 阶行列式称为 $A$ 的一个 $k$ 阶
\textbf{子式}。若 $A$ 中存在一个 $r$ 阶子式不为零，而所有 $r+1$ 阶子式（若存在）
全为零，则称 $A$ 的\textbf{秩}为 $r$，记作 $\rank(A)=r$。零矩阵的秩规定为
$0$。
\end{definition}

按定义，一个矩阵的秩就是其最高阶非零子式的阶数。由行列式性质
（性质~\ref{prop:det-transpose}：转置不变），$A$ 的子式与 $A\trans$ 的子式
互为转置，故
\[
  \rank(A\trans)=\rank(A).
\]
$n$ 阶方阵 $A$ 满秩（即 $\rank(A)=n$）当且仅当 $|A|\neq 0$，即 $A$ 可逆。

用子式直接求秩需要计算大量行列式，实际计算利用下面的定理。

\begin{theorem}\label{thm:rank-invariant}
初等变换不改变矩阵的秩。
\end{theorem}
\begin{proof}
由定理~\ref{thm:elem-mult}，只需证初等变换前后“存在非零 $s$ 阶子式”这一性质
不变。前两类变换（对换、数乘）显然保持。对第三类变换 $r_i+kr_j$，任取变换后
矩阵的一个 $s$ 阶子式 $D'$：若其行不涉及第 $i$ 行，则 $D'$ 是原矩阵的同位
子式；若涉及第 $i$ 行，由行列式性质~\ref{prop:det-split} 与
\ref{prop:det-equal-rows}，$D'$ 或者等于原矩阵的一个 $s$ 阶子式（当所选行
同时含第 $i,j$ 行时，拆出的第二项有两行相同而为零），或者等于原矩阵某个
$s$ 阶子式与 $k$ 乘另一个 $s$ 阶子式之和。总之，变换后的非零 $s$ 阶子式都
由原矩阵的非零 $s$ 阶子式产生；又因该变换可逆（用 $r_i-kr_j$ 还原），反向
亦然。故最高阶非零子式的阶数不变。
\end{proof}

\begin{theorem}\label{thm:rank-echelon}
行阶梯形矩阵的秩等于其非零行的个数。
\end{theorem}
\begin{proof}
设行阶梯形矩阵有 $r$ 个非零行，第 $i$ 行主元位于第 $j_i$ 列，则
$j_1<j_2<\cdots<j_r$。取这 $r$ 行与第 $j_1,\cdots,j_r$ 列，所得 $r$ 阶子式为
上三角行列式，主对角线元素为各主元，故不为零（例~\ref{ex:upper-tri}）。
而任何 $r+1$ 阶子式必含一个零行，故为零。因此秩为 $r$。
\end{proof}

于是求秩的方法是：用初等行变换把矩阵化为行阶梯形，非零行数即秩。

\begin{example}
求 $A=\begin{pmatrix} 1 & 2 & 3 & 4 \\ 2 & 4 & 6 & 9 \\ 1 & 1 & 1 & 1 \end{pmatrix}$
的秩。
\end{example}
\begin{proof}[解]
\[
  A\xrightarrow[\ r_3-r_1\ ]{r_2-2r_1}
  \begin{pmatrix} 1 & 2 & 3 & 4 \\ 0 & 0 & 0 & 1 \\ 0 & -1 & -2 & -3 \end{pmatrix}
  \xrightarrow{\ r_2\leftrightarrow r_3\ }
  \begin{pmatrix} 1 & 2 & 3 & 4 \\ 0 & -1 & -2 & -3 \\ 0 & 0 & 0 & 1 \end{pmatrix},
\]
行阶梯形有 $3$ 个非零行，故 $\rank(A)=3$。
\end{proof}

\begin{theorem}\label{thm:rank-product}
\begin{enumerate}[label=(\arabic*), itemsep=0pt]
  \item 若 $P$、$Q$ 可逆，则 $\rank(PAQ)=\rank(A)$；
  \item $\rank(AB)\le\min\{\rank(A),\rank(B)\}$；
  \item $\rank(A+B)\le\rank(A)+\rank(B)$（证明见第三章习题）。
\end{enumerate}
\end{theorem}
\begin{proof}
（1）由定理~\ref{thm:elem-product}，$P$、$Q$ 都是初等矩阵之积，而初等变换
不改变秩（定理~\ref{thm:rank-invariant}），故 $\rank(PAQ)=\rank(A)$。

（2）任何 $m\times n$ 矩阵 $B$ 都可经初等变换化为等价标准形
$B\sim\begin{pmatrix} E_r & O \\ O & O \end{pmatrix}$（对定理~\ref{thm:echelon}
得到的行最简形再作列变换：把主元列换到前 $r$ 列），其中 $r=\rank(B)$。即存在
可逆矩阵 $P,Q$ 使 $B=P\begin{pmatrix} E_r & O \\ O & O \end{pmatrix}Q$。于是
\[
  AB=(AP)\begin{pmatrix} E_r & O \\ O & O \end{pmatrix}Q,
\]
由（1），$\rank(AB)$ 等于 $AP\begin{pmatrix} E_r & O \\ O & O \end{pmatrix}$
的秩，而后者最后 $n-r$ 列全为零，其非零子式只能落在前 $r$ 列中，故秩 $\le r$。
于是 $\rank(AB)\le\rank(B)$。又
$\rank(AB)=\rank\big((AB)\trans\big)=\rank(B\trans A\trans)\le\rank(A\trans)=\rank(A)$。
\end{proof}

\begin{insight}[秩的几何意义：变换后还剩几维]
$r(A)$ 恰好是变换 $A$ 的\textbf{像空间的维数}——全体输出向量 $Ax$ 张成的空间
的维数（也就是列向量组的秩，见第三章定理~\ref{thm:col-rank}）。可逆时
$r(A)=n$：整个空间都被保留；$r(A)<n$ 时空间被压到一个 $r$ 维的“薄层”上。
例如 $\operatorname{diag}(1,0,0)$ 把三维空间压成一条直线，秩为 $1$。由此，
“初等变换不改变秩”（定理~\ref{thm:rank-invariant}）的几何内容就是：左乘
可逆矩阵只是一些扭转、拉伸与剪切的复合，不会降低空间的维数。
\end{insight}

\begin{qa}{求秩时为什么教材都强调“只用行变换”？用列变换或行列混用行吗？}
答：行变换保持\textbf{列向量之间的线性关系}（第三章引理~\ref{lem:col-relations}：
$Ak=0\iff(PA)k=0$），因此化为阶梯形后，既能数出秩，又能顺便读出列之间的
表示关系——这是第三章求极大无关组的标准做法。列变换会打乱这种关系（相当于
给未知量换名），但若\textbf{只求秩}，行、列变换都可以用（定理
~\ref{thm:rank-invariant} 对列同样成立）。习惯上只用行变换，是为了“一次
消元，同时得到秩与向量组的关系”。
\end{qa}

\chapterexercises{二}

\begin{enumerate}[label=\arabic*., itemsep=2pt]
\item 设
  $A=\begin{pmatrix} 1 & 2 \\ 3 & 4 \end{pmatrix}$，
  $B=\begin{pmatrix} 0 & 1 \\ 1 & 0 \end{pmatrix}$，求 $AB$、$BA$、$A^2-3A+2E$。
\item 计算下列乘积：
  $(1\ 2\ 3)\begin{pmatrix} 4 \\ 5 \\ 6 \end{pmatrix}$ 与
  $\begin{pmatrix} 4 \\ 5 \\ 6 \end{pmatrix}(1\ 2\ 3)$。
\item 求下列矩阵的逆：
  \[
  \text{(1)}\ \begin{pmatrix} 3 & 1 \\ 5 & 2 \end{pmatrix};
  \qquad
  \text{(2)}\ \begin{pmatrix} 1 & 2 & 3 \\ 2 & 2 & 3 \\ 1 & 1 & 1 \end{pmatrix}.
  \]
\item 设 $A=\begin{pmatrix} 1 & -1 \\ 2 & 3 \end{pmatrix}$，求
  $A^2-2A+5E$，并利用 $A^2-2A+5E=O$ 证明 $A$ 可逆并求 $A^{-1}$。
\item 解矩阵方程
  $\begin{pmatrix} 2 & 1 \\ 1 & 1 \end{pmatrix}X
   =\begin{pmatrix} 1 & 2 \\ 3 & 4 \end{pmatrix}$。
\item 设 $A$、$B$ 都是 $n$ 阶对称矩阵，证明 $AB$ 是对称矩阵的充分必要条件是
  $AB=BA$。
\item 证明：任一 $n$ 阶方阵都可以表示为一个对称矩阵与一个反对称矩阵之和。
\item 利用分块求
  $\begin{pmatrix} 5 & 6 & 0 & 0 \\ 6 & 7 & 0 & 0 \\ 0 & 0 & 1 & 2 \\ 0 & 0 & 3 & 5 \end{pmatrix}$
  的逆矩阵。
\item 用初等行变换求
  $\begin{pmatrix} 1 & 2 & 3 \\ 2 & 2 & 3 \\ 1 & 1 & 1 \end{pmatrix}$ 的逆
  （与例~\ref{ex:inv-3} 的结果一致）。
\item 求下列矩阵的秩：
  \[
  \text{(1)}\ \begin{pmatrix} 1 & 2 & 3 \\ 2 & 4 & 6 \\ 1 & 1 & 1 \end{pmatrix};
  \qquad
  \text{(2)}\ \begin{pmatrix} 1 & 2 & 3 & 4 \\ 2 & 4 & 6 & 9 \\ 3 & 6 & 9 & 13 \end{pmatrix}.
  \]
\item 设 $A$ 为 $n$ 阶可逆矩阵，证明 $|A^*|=|A|^{n-1}$。
\item 举例说明：存在矩阵 $A\neq O$、$B\neq O$ 使 $AB=O$；并由此说明消去律
  一般不成立。
\end{enumerate}

\chapter{向量组的线性相关性}

有了矩阵作为工具，本章研究 $n$ 维向量以及线性相关性理论。核心概念是线性组合、
线性相关与线性无关、极大无关组与向量组的秩；章末介绍向量空间的初步知识。

\section{$n$ 维向量}

\begin{definition}\label{def:n-vector}
$n$ 个有序的实数
\[
  \alpha=(a_1,a_2,\cdots,a_n)
\]
称为一个\textbf{$n$ 维向量}，其中 $a_i$ 称为向量 $\alpha$ 的第 $i$ 个
\textbf{分量}。向量也可写成列的形式；$n$ 维列向量就是 $n\times1$ 矩阵，
$n$ 维行向量就是 $1\times n$ 矩阵。全体 $n$ 维实向量构成的集合记作 $\R^n$。
\end{definition}

设 $\alpha=(a_1,\cdots,a_n)$，$\beta=(b_1,\cdots,b_n)$，$k$ 为数，定义
\[
  \alpha+\beta=(a_1+b_1,\cdots,a_n+b_n),\qquad
  k\alpha=(ka_1,\cdots,ka_n),\qquad
  -\alpha=(-1)\alpha .
\]
上述运算满足通常的八条运算律（加法交换律、结合律，零向量 $0=(0,\cdots,0)$
的存在性，负向量的存在性，以及数乘的四条分配律）。两个向量当且仅当对应分量
都相等时相等。

在几何上，$\R^2$ 与 $\R^3$ 中的向量就是平面与空间中的有向线段（或点），分量
就是坐标。$n>3$ 时虽无直观图像，代数运算完全相同。

\section{线性组合}

\begin{definition}\label{def:combination}
设 $\alpha_1,\cdots,\alpha_m$ 是一组向量。形如
\[
  k_1\alpha_1+k_2\alpha_2+\cdots+k_m\alpha_m
\]
的向量称为向量组 $\alpha_1,\cdots,\alpha_m$ 的一个\textbf{线性组合}（$k_i$
称为\textbf{系数}），也称该向量可由 $\alpha_1,\cdots,\alpha_m$ \textbf{线性
表示}。$\alpha_1,\cdots,\alpha_m$ 的全体线性组合构成的集合记作
$\operatorname{span}(\alpha_1,\cdots,\alpha_m)$。
\end{definition}

例如，任一向量 $\alpha=(a_1,a_2,a_3)$ 都是标准单位向量
$\varepsilon_1=(1,0,0)$、$\varepsilon_2=(0,1,0)$、$\varepsilon_3=(0,0,1)$
的线性组合：$\alpha=a_1\varepsilon_1+a_2\varepsilon_2+a_3\varepsilon_3$。

一个向量能否由某向量组线性表示，本质上是一个线性方程组是否有解的问题：把各
向量写成列向量，$\beta=\sum k_i\alpha_i$ 即 $Ak=\beta$，其中
$A=(\alpha_1,\cdots,\alpha_m)$。其结论由秩刻画。

\begin{theorem}\label{thm:represent-rank}
$\beta$ 可由 $\alpha_1,\cdots,\alpha_m$ 线性表示的充分必要条件是
\[
  \rank(\alpha_1,\cdots,\alpha_m,\beta)=\rank(\alpha_1,\cdots,\alpha_m);
\]
且表示式唯一的充分必要条件是 $\alpha_1,\cdots,\alpha_m$ 线性无关
（见 \S\ref{sec:dependence}）。
\end{theorem}

该定理在 \S\ref{sec:vec-rank} 建立向量组的秩的理论后再证明。

\begin{example}
$\beta=(1,2,3)$ 能否由 $\alpha_1=(1,1,0)$、$\alpha_2=(0,1,1)$ 线性表示？
\end{example}
\begin{proof}[解]
设 $k_1\alpha_1+k_2\alpha_2=\beta$，比较分量得
\[
  k_1=1,\qquad k_1+k_2=2,\qquad k_2=3,
\]
此方程组无解（前两式给出 $k_2=1$，第三式给出 $k_2=3$，矛盾）。故 $\beta$
不能由 $\alpha_1,\alpha_2$ 线性表示。
\end{proof}

\section{线性相关与线性无关}\label{sec:dependence}

\begin{definition}\label{def:dependence}
若存在一组\textbf{不全为零}的数 $k_1,\cdots,k_m$，使
\[
  k_1\alpha_1+k_2\alpha_2+\cdots+k_m\alpha_m=0 ,
\]
则称向量组 $\alpha_1,\cdots,\alpha_m$ \textbf{线性相关}；否则称它
\textbf{线性无关}：上式成立当且仅当 $k_1=\cdots=k_m=0$。
\end{definition}

等价地说：把 $\alpha_1,\cdots,\alpha_m$ 作为列排成矩阵 $A$，向量组线性相关
当且仅当齐次线性方程组 $Ak=0$ 有非零解。

几何上：两个向量线性相关当且仅当其中一个是另一个的倍数（共线）；$\R^3$ 中
三个向量线性相关当且仅当它们共面（在同一张过原点的平面上）。

\begin{theorem}\label{thm:dependent-one-of}
向量组 $\alpha_1,\cdots,\alpha_m$ 线性相关的充分必要条件是其中至少有一个
向量可由其余向量线性表示。
\end{theorem}
\begin{proof}
若 $\sum k_i\alpha_i=0$ 且某个 $k_j\neq0$，则可解出
$\alpha_j=-\frac{1}{k_j}\sum_{i\neq j}k_i\alpha_i$。反之，若
$\alpha_j=\sum_{i\neq j}l_i\alpha_i$，则
$\alpha_j-\sum_{i\neq j}l_i\alpha_i=0$ 就是一个非平凡的表示式。
\end{proof}

\begin{theorem}\label{thm:unique-representation}
若向量组 $\alpha_1,\cdots,\alpha_m$ 线性无关，而 $\alpha_1,\cdots,\alpha_m,\beta$
线性相关，则 $\beta$ 可由 $\alpha_1,\cdots,\alpha_m$ 线性表示，且表示式唯一。
\end{theorem}
\begin{proof}
由相关性，存在不全为零的数使 $k_1\alpha_1+\cdots+k_m\alpha_m+k\beta=0$。若
$k=0$，则 $\sum k_i\alpha_i=0$ 是非平凡的关系，与无关性矛盾；故 $k\neq0$，
于是 $\beta=-\frac{1}{k}\sum k_i\alpha_i$。唯一性：若又有
$\beta=\sum l_i\alpha_i$，两式相减得 $\sum(k_i-l_i)\alpha_i=0$，由无关性
$k_i=l_i$。
\end{proof}

\begin{theorem}\label{thm:subset}
若向量组的一个部分组线性相关，则整个向量组线性相关；等价地，线性无关向量组
的任何部分组都线性无关。
\end{theorem}
\begin{proof}
把部分组上的非平凡关系在其余向量上取零系数扩充，即得整组的一个非平凡关系。
\end{proof}

\begin{theorem}\label{thm:more-than-n}
$\R^n$ 中任意 $m>n$ 个向量必线性相关。
\end{theorem}
\begin{proof}
设 $A=(\alpha_1,\cdots,\alpha_m)$，这是 $n\times m$ 矩阵。把 $A$ 用初等行
变换化为行阶梯形，其主元个数不超过行数 $n$，故主元列至多 $n$ 个；而 $m>n$，
必有某列不是主元列。在阶梯形方程组 $Ak=0$ 中回代求解：自由未知量取非零值，
即得非零解 $k$（自最下面的主元行逐行回代，解出各主元未知量后，自由未知量
可任取非零值）。
\end{proof}

\begin{example}
判断 $\alpha_1=(1,1,0)$、$\alpha_2=(0,1,1)$、$\alpha_3=(1,0,-1)$ 是否线性
相关。
\end{example}
\begin{proof}[解]
解 $k_1\alpha_1+k_2\alpha_2+k_3\alpha_3=0$，按分量写出：
\[
  \begin{cases}
    k_1\hphantom{+k_2}+k_3=0, \\
    k_1+k_2\hphantom{+k_3}=0, \\
    \hphantom{k_1+}k_2-k_3=0.
  \end{cases}
\]
前两式给出 $k_2=k_3=-k_1$，代入第三式得 $-k_1+k_1=0$ 恒成立，故 $k_1$ 可
任取。取 $k_1=1$ 得非平凡关系 $\alpha_1-\alpha_2-\alpha_3=0$，向量组线性
相关。
\end{proof}

\begin{example}
$n$ 个标准单位向量 $\varepsilon_1,\cdots,\varepsilon_n$ 线性无关：由
$\sum k_i\varepsilon_i=(k_1,\cdots,k_n)=0$ 立得每个 $k_i=0$。
\end{example}

\section{向量组的秩}\label{sec:vec-rank}

\begin{definition}\label{def:maximal-independent}
设 $S$ 是一个向量组。若 $S$ 的一个部分组 $\alpha_{i_1},\cdots,\alpha_{i_r}$
线性无关，且 $S$ 中任一向量都可由它线性表示，则称该部分组是 $S$ 的一个
\textbf{极大线性无关组}（简称\textbf{极大无关组}），数 $r$ 称为向量组 $S$
的\textbf{秩}。
\end{definition}

按定义，秩似乎依赖于所选的极大无关组；下面的结果表明它是唯一确定的。先给出
一个联系列向量组关系与初等行变换的引理。

\begin{lemma}\label{lem:col-relations}
设 $P$ 可逆，则矩阵 $A$ 的列向量组与 $PA$ 的列向量组满足完全相同的线性关系，
即 $Ak=0\iff(PA)k=0$。
\end{lemma}
\begin{proof}
$Ak=0\Rightarrow(PA)k=P(Ak)=0$；反之 $(PA)k=0\Rightarrow
Ak=P^{-1}\big((PA)k\big)=0$。
\end{proof}

\begin{theorem}\label{thm:col-rank}
设 $A$ 经初等行变换化为有 $r$ 个主元的行最简形 $R$，则：
（1）$A$ 的 $r$ 个主元列线性无关，$A$ 的其余列都可由主元列线性表示；
（2）$A$ 的列向量组的秩（列秩）等于 $r=\rank(A)$；
（3）$A$ 的行向量组的秩（行秩）等于列秩。
\end{theorem}
\begin{proof}
在 $R$ 中，各主元列恰为单位向量 $e_1,\cdots,e_r$，它们线性无关；每个非主元
列以自身各元素为系数可由这些单位向量线性表示。由引理~\ref{lem:col-relations}，
$A$ 的列之间也满足相同的关系，故（1）成立，进而（2）由定理~\ref{thm:rank-echelon}
（秩 $=r$）得证。（3）对 $A\trans$ 用（2）：$A\trans$ 的列秩等于
$\rank(A\trans)=\rank(A)$，而 $A\trans$ 的列正是 $A$ 的行。
\end{proof}

由该定理，向量组的秩等于以这些向量为列排成矩阵的秩，并可由行阶梯形直接读出；
\textbf{原矩阵的主元列}就构成向量组的一个极大无关组。

\begin{theorem}\label{thm:represent-inequality}
若线性无关向量组 $A$ 含 $r$ 个向量，且可由含 $s$ 个向量的向量组 $B$ 线性
表示，则 $r\le s$。因此，同一向量组的任意两个极大无关组所含向量个数相同，
秩是唯一确定的。
\end{theorem}
\begin{proof}
把两组向量按列排成矩阵 $A=(\alpha_1,\cdots,\alpha_r)$、
$B=(\beta_1,\cdots,\beta_s)$。“$A$ 组可由 $B$ 组表示”即存在 $s\times r$
矩阵 $K$ 使 $A=BK$。由定理~\ref{thm:rank-product}，
$\rank(A)\le\rank(B)\le s$；又 $A$ 的列线性无关，由定理~\ref{thm:col-rank}，
$\rank(A)=r$。故 $r\le s$。若两个极大无关组分别含 $r$、$s$ 个向量，则它们
互相可表示（都表示整个向量组），于是 $r\le s$ 且 $s\le r$，即 $r=s$。
\end{proof}

\begin{proof}[定理~\ref{thm:represent-rank} 的证明]
$\beta$ 可由 $\alpha_1,\cdots,\alpha_m$ 表示，等价于在 $(\alpha_1,\cdots,
\alpha_m)$ 后添列 $\beta$ 秩不增大，即两组秩相等（添列只可能使秩不减）。
唯一性即定理~\ref{thm:unique-representation}。
\end{proof}

\begin{example}
求向量组 $\alpha_1=(1,0,2)$、$\alpha_2=(0,1,1)$、$\alpha_3=(1,1,3)$、
$\alpha_4=(2,1,5)$ 的一个极大无关组与秩，并把其余向量用该极大无关组线性
表示。
\end{example}
\begin{proof}[解]
把各向量按列排成矩阵并作初等行变换：
\[
  (\alpha_1,\alpha_2,\alpha_3,\alpha_4)=
  \begin{pmatrix} 1 & 0 & 1 & 2 \\ 0 & 1 & 1 & 1 \\ 2 & 1 & 3 & 5 \end{pmatrix}
  \xrightarrow{\ r_3-2r_1\ }
  \begin{pmatrix} 1 & 0 & 1 & 2 \\ 0 & 1 & 1 & 1 \\ 0 & 1 & 1 & 1 \end{pmatrix}
  \xrightarrow{\ r_3-r_2\ }
  \begin{pmatrix} 1 & 0 & 1 & 2 \\ 0 & 1 & 1 & 1 \\ 0 & 0 & 0 & 0 \end{pmatrix}.
\]
主元在第 1、2 列：$\alpha_1,\alpha_2$ 是一个极大无关组，向量组的秩为 $2$；
由非主元列读出
\[
  \alpha_3=\alpha_1+\alpha_2,\qquad \alpha_4=2\alpha_1+\alpha_2 .
\]
\end{proof}

\section{向量空间、基与坐标}

\begin{definition}\label{def:vector-space}
设 $V$ 是若干个同维向量构成的非空集合。若 $V$ 对向量的加法与数乘封闭
（$\alpha,\beta\in V$ 时 $\alpha+\beta\in V$ 且 $k\alpha\in V$），则称 $V$
是一个\textbf{向量空间}（$\R^n$ 的线性子空间）。任何这样的 $V$ 都含有零向量。
\end{definition}

例如：$\R^n$ 本身；$\R^3$ 中过原点的直线或平面；齐次线性方程组的解集
（见第四章）；以及 $\operatorname{span}(\alpha_1,\cdots,\alpha_m)$——它是
包含 $\alpha_1,\cdots,\alpha_m$ 的最小向量空间。

\begin{definition}\label{def:basis}
若向量空间 $V$ 中的向量组 $\beta_1,\cdots,\beta_n$ 线性无关，且 $V$ 中任一
向量都可由它们线性表示，则称 $\beta_1,\cdots,\beta_n$ 是 $V$ 的一组\textbf{基}。
此时每个 $\alpha\in V$ 的表示式 $\alpha=x_1\beta_1+\cdots+x_n\beta_n$ 唯一
（定理~\ref{thm:unique-representation}），系数 $(x_1,\cdots,x_n)$ 称为
$\alpha$ 在该基下的\textbf{坐标}，$n$ 称为 $V$ 的\textbf{维数}（由定理
~\ref{thm:represent-inequality}，任何两组基所含向量个数相同）。
\end{definition}

$\R^n$ 的标准基是 $\varepsilon_1,\cdots,\varepsilon_n$。一个常用判别法：
$\R^n$ 中 $n$ 个向量构成基的充分必要条件是以它们为列的矩阵的行列式非零
（由定理~\ref{thm:more-than-n} 与第一章：$n$ 个 $n$ 维向量线性无关当且仅当
行列式非零）。

\begin{example}
求 $\beta=(1,2,1)$ 在基 $\alpha_1=(1,1,0)$、$\alpha_2=(0,1,1)$、
$\alpha_3=(1,0,1)$ 下的坐标。
\end{example}
\begin{proof}[解]
解 $k_1\alpha_1+k_2\alpha_2+k_3\alpha_3=\beta$：
\[
  \begin{cases}
    k_1\hphantom{+k_2}+k_3=1, \\
    k_1+k_2\hphantom{+k_3}=2, \\
    \hphantom{k_1+}k_2+k_3=1.
  \end{cases}
\]
三式相加得 $2(k_1+k_2+k_3)=4$，即 $k_1+k_2+k_3=2$；用它分别减去三式得
$k_2=1$、$k_3=1$、$k_1=1$。检验：$\alpha_1+\alpha_2=(1,2,1)=\beta$。故坐标
为 $(1,1,0)$。
\end{proof}

设 $\alpha_1,\cdots,\alpha_n$ 与 $\beta_1,\cdots,\beta_n$ 是 $V$ 的两组基，
且 $(\beta_1,\cdots,\beta_n)=(\alpha_1,\cdots,\alpha_n)C$（$C$ 的第 $j$ 列是
$\beta_j$ 在 $\alpha$ 基下的坐标），则 $C$ 可逆，且坐标变换公式为
$x_\alpha=C\,x_\beta$；$C$ 称为由 $\alpha$ 基到 $\beta$ 基的\textbf{过渡矩阵}。

\chapterexercises{三}

\begin{enumerate}[label=\arabic*., itemsep=2pt]
\item 判断 $(1,1,1)$、$(1,2,3)$、$(1,3,6)$ 是否线性无关。
\item $\beta=(1,2,1)$ 能否由 $\alpha_1=(1,1,0)$、$\alpha_2=(0,1,1)$ 线性
  表示？若能，表示式是否唯一？
\item 求向量组 $(1,2,3)$、$(2,4,6)$、$(0,1,2)$ 的一个极大无关组与秩。
\item 求向量组 $(1,0,2)$、$(0,1,1)$、$(1,1,3)$、$(2,1,5)$ 的秩、一个极大
  无关组，并用它表示其余向量。
\item 证明定理~\ref{thm:subset}：线性无关向量组的任何部分组都线性无关。
\item 证明：含零向量的向量组必线性相关。
\item 证明：两个向量线性相关的充分必要条件是它们的对应分量成比例。
\item 通过求解相应的齐次线性方程组，证明 $\R^n$ 中任何 $n+1$ 个向量必线性
  相关（定理~\ref{thm:more-than-n} 取 $m=n+1$）。
\item $\lambda$ 取何值时，$(1,1,0)$、$(0,1,1)$、$(\lambda,0,1)$ 线性相关？
  对该 $\lambda$ 写出一个非平凡的线性关系式。
\item 求 $(2,3,3)$ 在基 $(1,1,0)$、$(0,1,1)$、$(1,0,1)$ 下的坐标。
\item 证明 $\operatorname{span}(\alpha_1,\cdots,\alpha_m)$ 是向量空间，且是
  包含 $\alpha_1,\cdots,\alpha_m$ 的最小向量空间。
\item 证明 $\rank(A+B)\le\rank(A)+\rank(B)$。（提示：先证
  $\rank(A\ \ B)\le\rank(A)+\rank(B)$，其中 $(A\ \ B)$ 表示把 $B$ 的列接在
  $A$ 的右边；再利用 $A+B=(A\ \ B)\begin{pmatrix} E \\ E \end{pmatrix}$。）
\end{enumerate}

% Stepwise build: re-enable each include as the chapter is written.
% \chapter{线性方程组}

第一章的克拉默法则只解决了最特殊的情形：方程个数与未知数个数相等，且系数
行列式非零。本章用秩作工具，把\textbf{任意}线性方程组的三大问题一次说清：
什么时候有解？解是否唯一？解不唯一时，全部解长什么样？

\begin{insight}[观点：方程 = 约束，解方程组 = 求交集]
一个线性方程就是对位置的一条规定：二元方程画出平面上的一条\textbf{直线}，
三元方程画出空间里的一张\textbf{平面}。解方程组，就是求这些直线、平面的
\textbf{公共点}。三条直线可能交于一点、交成一片，也可能永远碰不到一起——
“有没有解、有多少解”因此有了画面；本章的判定定理（秩判别法）就是把这个
画面翻译成代数。
\end{insight}

\section{消元法与解的判定定理}\label{sec:elimination}

设有 $m$ 个方程、$n$ 个未知数的线性方程组 $Ax=b$。把系数与常数并排写成
\[
  B=(A\mid b),
\]
称为\textbf{增广矩阵}。求解的标准动作是：对 $B$ 作初等行变换（消元的机械化
形式，第二章 \S\ref{sec:elementary}），化为行阶梯形——每作一次行变换，方程组
同解。

\begin{theorem}[解的判定定理]\label{thm:solvability}
设 $A$ 为 $m\times n$ 系数矩阵，$B=(A\mid b)$ 为增广矩阵，则：
\begin{enumerate}[label=(\arabic*), itemsep=0pt]
  \item $Ax=b$ 有解 $\iff$ $\rank(A)=\rank(B)$；
  \item 有解时，若 $\rank(A)=n$（未知数个数），解唯一；若 $\rank(A)=r<n$，
        有无穷多解。
\end{enumerate}
\end{theorem}
\begin{proof}
对 $B$ 作行变换化为行阶梯形，方程组同解。阶梯形中出现矛盾行
$(0,\cdots,0\mid d)$（$d\neq0$，即 $0=d$）当且仅当 $B$ 的非零行比 $A$ 的多，
即 $\rank(B)=\rank(A)+1$；此时无解。若无矛盾行，则 $\rank(B)=\rank(A)$，
回代可解出全部主元未知数；自由未知量（非主元列对应的未知数）的个数为
$n-r$。$r=n$ 时没有自由未知量，解唯一；$r<n$ 时自由未知量可任取，解有无穷
多个。
\end{proof}

对 $n$ 个未知数、$n$ 个方程的情形，$\rank(A)=n\iff|A|\neq0$，判定定理与
克拉默法则的结论一致——秩判别法是克拉默法则的推广，且不再要求方阵。

\begin{figure}[htbp]
\centering
\begin{tikzpicture}[scale=0.78, >=Stealth]
  % (a) unique solution: three planes through one point
  \begin{scope}
    \draw[fill=blue!10, draw=blue!50!black, opacity=0.9]
      (0.1,0.55) -- (1.55,0.1) -- (1.85,1.05) -- (0.4,1.5) -- cycle;
    \draw[fill=red!10, draw=red!50!black, opacity=0.9]
      (0.35,0.2) -- (1.9,0.55) -- (1.55,1.55) -- (0.0,1.2) -- cycle;
    \draw[fill=green!10, draw=green!45!black, opacity=0.9]
      (0.5,0.0) -- (1.95,0.35) -- (1.75,1.35) -- (0.3,1.0) -- cycle;
    \fill (1.18,0.83) circle (1.9pt);
    \node[below] at (1.1,-0.55) {(a) $\rank(A)=\rank(B)=3$：唯一解};
  \end{scope}
  % (b) a common line
  \begin{scope}[xshift=4.9cm]
    \draw[very thick, black!70] (0.15,0.75) -- (2.0,0.45);
    \draw[fill=blue!10, draw=blue!50!black, opacity=0.9]
      (0.0,0.6) -- (1.7,0.15) -- (1.95,1.0) -- (0.25,1.45) -- cycle;
    \draw[fill=red!10, draw=red!50!black, opacity=0.9]
      (0.25,0.35) -- (1.95,0.6) -- (1.7,1.6) -- (0.0,1.35) -- cycle;
    \draw[fill=green!10, draw=green!45!black, opacity=0.9]
      (0.3,0.85) -- (2.05,0.5) -- (2.2,1.35) -- (0.45,1.7) -- cycle;
    \node[below] at (1.1,-0.55) {(b) $\rank(A)=\rank(B)=2<3$：无穷多解};
  \end{scope}
  % (c) no common point: two parallel planes
  \begin{scope}[xshift=9.8cm]
    \draw[fill=blue!10, draw=blue!50!black, opacity=0.9]
      (0.0,0.2) -- (1.6,0.0) -- (1.9,0.9) -- (0.3,1.1) -- cycle;
    \draw[fill=blue!25, draw=blue!50!black, opacity=0.9]
      (0.45,0.75) -- (2.05,0.55) -- (2.35,1.45) -- (0.75,1.65) -- cycle;
    \node at (1.15,2.0) {\small 平行，永不相交};
    \node[below] at (1.1,-0.55) {(c) $\rank(A)<\rank(B)$：无解};
  \end{scope}
\end{tikzpicture}
\caption{三元方程组的解集 = 三张平面的公共点，三种命运由秩完全决定。}
\label{fig:three-planes}
\end{figure}

\begin{qa}{有了行列式判别（$|A|\neq0$），为什么还要引入秩来判别？}
答：行列式只对\textbf{方阵}才有意义，而实际方程组常常“方程数 $\neq$ 未知数
数”，或系数行列式恰为零——这些情形行列式根本无从开口。秩对任意矩阵都有
定义，一次覆盖所有情形：比较 $r(A)$ 与 $r(A\mid b)$ 定有解无解，比较 $r(A)$
与 $n$ 定唯一无穷。可以说：行列式回答“这张 $n\times n$ 表满不满秩”，秩回答
“任意一张表还剩几维信息”。
\end{qa}

\begin{example}\label{ex:param-system}
讨论 $\lambda$ 取何值时，方程组
\[
  \begin{cases}
    \lambda x_1+x_2+x_3=1, \\
    x_1+\lambda x_2+x_3=\lambda, \\
    x_1+x_2+\lambda x_3=\lambda^2
  \end{cases}
\]
无解、有唯一解、有无穷多解。
\end{example}
\begin{proof}[解]
先算系数行列式。把各列都加到第 1 列，第 1 列各元素均为 $\lambda+2$，提出
公因子后再用 $r_i-r_1$ 化简，得
\[
  |A|=\begin{vmatrix} \lambda & 1 & 1 \\ 1 & \lambda & 1 \\ 1 & 1 & \lambda \end{vmatrix}
      =(\lambda+2)\begin{vmatrix} 1 & 1 & 1 \\ 1 & \lambda & 1 \\ 1 & 1 & \lambda \end{vmatrix}
      =(\lambda+2)(\lambda-1)^2 .
\]
（1）当 $\lambda\neq1$ 且 $\lambda\neq-2$ 时 $|A|\neq0$，方程组有唯一解。

（2）当 $\lambda=1$ 时，三个方程都是 $x_1+x_2+x_3=1$：$r(A)=r(B)=1<n$，
有无穷多解。

（3）当 $\lambda=-2$ 时，三式相加：左端每列系数之和为 $0$，右端为
$1-2+4=3\neq0$，出现矛盾方程 $0=3$，故 $r(B)=3>2=r(A)$，无解。
\end{proof}

\section{齐次线性方程组}\label{sec:homogeneous}

齐次方程组 $Ax=0$ 总有零解，因此问题只有一个：\textbf{何时有非零解，非零解
有多少、长什么样？}

\begin{proposition}\label{prop:solution-space}
若 $\xi_1,\xi_2$ 都是 $Ax=0$ 的解，则 $\xi_1+\xi_2$ 与 $k\xi_1$ 也都是它的解。
因此 $Ax=0$ 的解集是一个向量空间（第三章定义~\ref{def:vector-space}），
称为\textbf{解空间}。
\end{proposition}
\begin{proof}
$A(\xi_1+\xi_2)=A\xi_1+A\xi_2=0$；$A(k\xi_1)=kA\xi_1=0$。封闭性得证。
\end{proof}

\begin{definition}\label{def:fundamental}
解空间的一组基称为 $Ax=0$ 的\textbf{基础解系}；若 $\xi_1,\cdots,\xi_t$ 是
基础解系，则 $Ax=0$ 的全部解为
\[
  x=c_1\xi_1+c_2\xi_2+\cdots+c_t\xi_t\qquad(c_1,\cdots,c_t\text{ 为任意常数}).
\]
\end{definition}

\begin{theorem}\label{thm:fundamental}
若 $\rank(A)=r<n$，则 $Ax=0$ 的基础解系恰含 $n-r$ 个向量，即解空间是
$\R^n$ 中一个 $n-r$ 维的（过原点的）“平面”。
\end{theorem}
\begin{proof}
对 $A$ 作行变换化为行最简形，主元列 $r$ 个，其余 $n-r$ 列对应的未知量为
自由未知量。不失一般性设自由未知量为 $x_{r+1},\cdots,x_n$，则每个解由自由
未知量的取值完全确定：
\[
  x=(x_1,\cdots,x_r,x_{r+1},\cdots,x_n)\trans,
  \qquad x_i\ (i\le r)\text{ 由回代确定}.
\]
对每个 $j=1,\cdots,n-r$，取 $x_{r+j}=1$、其余自由未知量为 $0$，得解 $\xi_j$。
一方面，$\xi_1,\cdots,\xi_{n-r}$ 的第 $r+1,\cdots,n$ 个分量排成单位矩阵，
故它们线性无关（部分分量无关则整体无关）。另一方面，任一解 $x$ 以其自由
分量 $c_{r+1},\cdots,c_n$ 为系数满足
$x=\sum_{j}c_{r+j}\xi_j$（两边自由分量相同，主元分量由同一回代公式确定）。
故 $\xi_1,\cdots,\xi_{n-r}$ 是解空间的基，含 $n-r$ 个向量。
\end{proof}

\begin{insight}[本质：被压掉的信息以解的身份“回来”了]
把 $A$ 看作变换：像空间（列空间）是 $r$ 维，零空间（解空间）是 $n-r$ 维，
两者恰好把 $\R^n$ 分摊干净：
\[
  \underbrace{r}_{\text{保留的维数}}+\underbrace{(n-r)}_{\text{被压成零的维数}}=n .
\]
这就是\textbf{维数公式} $r(A)+\dim N(A)=n$ 的内容。矩阵秩越高、“个性”越强，
能压成零的方向就越少；满秩方阵（$r=n$）把零空间压到只剩零向量，方程组
$Ax=0$ 只有零解——与第一章推论~\ref{cor:homog} 再次呼应。
\end{insight}

\begin{qa}{自由未知量的个数为什么恰好是 $n-r$？基础解系为什么恰好 $n-r$ 个向量？}
答：行最简形有 $r$ 个主元，每个主元“钉死”一个未知量（它由其余未知量表出），
剩下 $n-r$ 个自由列各贡献一个可任意取值的未知量——自由度是 $n-r$。构造基础
解系时，每次令一个自由未知量取 $1$、其余取 $0$，得到一个解；$n-r$ 个解的自由
分量恰排成单位矩阵，保证它们线性无关；而任何解的组合系数就是它自己的自由
分量，保证它们张满整个解空间。于是“不多不少，恰好 $n-r$ 个”。
\end{qa}

\begin{example}\label{ex:fundamental}
求齐次方程组 $\begin{cases} x_1+x_2+x_3+x_4=0, \\ 2x_1+3x_2+4x_3+5x_4=0 \end{cases}$
的基础解系与通解。
\end{example}
\begin{proof}[解]
\[
  A=\begin{pmatrix} 1 & 1 & 1 & 1 \\ 2 & 3 & 4 & 5 \end{pmatrix}
  \xrightarrow{\ r_2-2r_1\ }
  \begin{pmatrix} 1 & 1 & 1 & 1 \\ 0 & 1 & 2 & 3 \end{pmatrix}
  \xrightarrow{\ r_1-r_2\ }
  \begin{pmatrix} 1 & 0 & -1 & -2 \\ 0 & 1 & 2 & 3 \end{pmatrix},
\]
$r=2$，$x_3,x_4$ 自由，且 $x_1=x_3+2x_4$，$x_2=-2x_3-3x_4$。分别取
$(x_3,x_4)=(1,0)$ 与 $(0,1)$：
\[
  \xi_1=(1,-2,1,0)\trans,\qquad \xi_2=(2,-3,0,1)\trans,
\]
即基础解系（$n-r=2$ 个），通解为 $x=c_1\xi_1+c_2\xi_2$。可验证
$A\xi_1=A\xi_2=0$。
\end{proof}

\section{非齐次线性方程组：解的结构}\label{sec:nonhomogeneous}

\begin{proposition}\label{prop:nonhomog-diff}
设 $\eta_1,\eta_2$ 是 $Ax=b$（$b\neq0$）的解，则 $\eta_1-\eta_2$ 是 $Ax=0$
的解；又设 $\eta^*$ 是 $Ax=b$ 的一个解（\textbf{特解}），$\xi$ 是 $Ax=0$ 的
解，则 $\eta^*+\xi$ 仍是 $Ax=b$ 的解。
\end{proposition}
\begin{proof}
$A(\eta_1-\eta_2)=A\eta_1-A\eta_2=b-b=0$；$A(\eta^*+\xi)=b+0=b$。
\end{proof}

\begin{theorem}[解的结构定理]\label{thm:nonhomog-structure}
设 $Ax=b$（$b\neq0$）有解，$\eta^*$ 是它的一个特解，$\xi_1,\cdots,\xi_{n-r}$
是 $Ax=0$ 的基础解系（$r=\rank(A)$），则其\textbf{全部解}为
\[
  x=\eta^*+c_1\xi_1+\cdots+c_{n-r}\xi_{n-r}\qquad(c_i\text{ 为任意常数}).
\]
\end{theorem}
\begin{proof}
一方面，$\eta^*+c_1\xi_1+\cdots$ 形如“特解 $+$ 齐次解”，由命题
~\ref{prop:nonhomog-diff} 都是 $Ax=b$ 的解。另一方面，设 $x$ 是任一解，则
$x-\eta^*$ 是 $Ax=0$ 的解，故可由基础解系线性表示：
$x-\eta^*=\sum c_i\xi_i$。两个方向合起来即得全部解。
\end{proof}

\begin{figure}[htbp]
\centering
\begin{tikzpicture}[>=Stealth, scale=1.0]
  \draw[->] (-2.4,0) -- (2.8,0) node[right] {$x_1$};
  \draw[->] (0,-1.3) -- (0,2.4) node[above] {$x_2$};
  \draw[dashed, gray, very thick] (-2.3,1.15) -- (2.6,-1.3)
        node[below left] {\small 解空间 $S{:}\ Ax=0$};
  \draw[very thick, red!70!black] (-2.3,2.35) -- (2.6,-0.1)
        node[below right] {\small 解集 $\eta^*+S$};
  \draw[->, very thick, blue!60!black] (0,0) -- (0.5,1.75)
        node[above right] {$\eta^*$};
  \draw[->, gray] (0.5,1.75) -- (1.6,1.2) node[right] {\small $\xi$（沿 $S$ 方向滑动）};
\end{tikzpicture}
\caption{非齐次方程组的解集 = 把齐次解空间整体平移到特解处。特解是“锚点”，
齐次解给出在“平面”内自由滑动的方向。}
\label{fig:translate}
\end{figure}

\begin{qa}{为什么“特解 $+$ 齐次通解”就是全部解？增广矩阵多了一列，秩为什么至多多 1？}
答：第一个问题分两个方向：任取“特解 $+$ 齐次解”，代入验证仍是解（命题
~\ref{prop:nonhomog-diff}）；反之任一解减去特解就是齐次解（差的验证），
从而必落在齐次通解里。几何上（图~\ref{fig:translate}）：解集是一张\textbf{不过
原点}、与解空间平行的“平面”，特解给出它的锚点，齐次基础解系给出平面内的
滑动方向。第二个问题：增广矩阵只比 $A$ 多一列，列空间至多多出一个方向，
故 $r(A)\le r(B)\le r(A)+1$——“有解无解”的分歧恰好发生在这 $+1$ 上。
\end{qa}

\begin{example}\label{ex:nonhomog}
求 $\begin{cases} x_1+x_2+x_3=3, \\ x_1+2x_2+3x_3=6 \end{cases}$ 的通解。
\end{example}
\begin{proof}[解]
\[
  B=\begin{pmatrix} 1 & 1 & 1 &\vline& 3 \\ 1 & 2 & 3 &\vline& 6 \end{pmatrix}
  \xrightarrow{\ r_2-r_1\ }
  \begin{pmatrix} 1 & 1 & 1 &\vline& 3 \\ 0 & 1 & 2 &\vline& 3 \end{pmatrix}
  \xrightarrow{\ r_1-r_2\ }
  \begin{pmatrix} 1 & 0 & -1 &\vline& 0 \\ 0 & 1 & 2 &\vline& 3 \end{pmatrix}.
\]
$r(A)=r(B)=2<3$，有无穷多解。取 $x_3=0$ 得特解 $\eta^*=(0,3,0)\trans$；
令 $x_3=1$ 得齐次解 $\xi=(1,-2,1)\trans$。通解为
\[
  x=(0,3,0)\trans+c\,(1,-2,1)\trans\qquad(c\text{ 为任意常数}).\qedhere
\]
\end{proof}

\chapterexercises{四}

\begin{enumerate}[label=\arabic*., itemsep=2pt]
\item 解方程组 $\begin{cases} x_1+x_2+x_3=2, \\ x_1+2x_2+2x_3=3, \\ x_2+x_3=1. \end{cases}$
\item 求齐次方程组
  $\begin{cases} x_1+2x_2+3x_3=0, \\ 2x_1+4x_2+6x_3=0, \\ 3x_1+6x_2+9x_3=0 \end{cases}$
  的基础解系与通解。
\item 求通解：$\begin{cases} x_1+x_2+x_3+x_4=4, \\ x_1+2x_2+3x_3+4x_4=10. \end{cases}$
\item $\lambda$ 取何值时，齐次方程组
  $\begin{cases} (\lambda+1)x_1+x_2+x_3=0, \\ x_1+(\lambda+1)x_2+x_3=0, \\
  x_1+x_2+(\lambda+1)x_3=0 \end{cases}$
  有非零解？求出此时的基础解系。
\item 证明命题~\ref{prop:solution-space}：齐次方程组的解集是向量空间。
\item 证明定理~\ref{thm:nonhomog-structure} 的两个核心断言：非齐次两解之差
  是齐次解；非齐次解集等于“特解 $+$ 齐次解集”。
\item 解方程组 $\begin{cases} x_1+2x_2+3x_3=6, \\ 2x_1+3x_2+x_3=6, \\
  x_1+x_2+x_3=3. \end{cases}$
\item 证明：方程个数少于未知数个数的齐次方程组必有非零解（用秩，并指出它与
  定理~\ref{thm:more-than-n} 的关系）。
\item 设 $Ax=b$ 有解且 $\rank(A)=r<n$，问通解中含有几个任意常数？为什么？
\item 证明：$\rank(A)\le\rank(A\mid b)\le\rank(A)+1$。（提示：增广矩阵比 $A$
  只多一列，列空间至多扩大一维。）
\end{enumerate}

% \chapter{特征值与特征向量}

本章研究方阵最内在的结构：哪些方向在变换下\textbf{保持不变}？答案是特征向量；
不变的倍数是特征值。作为准备，先介绍内积、长度与正交性；随后讨论相似矩阵
与对角化；最后证明实对称矩阵必可正交对角化——它是第六章化简二次型的钥匙。

\section{向量的内积、长度与正交性}

\begin{definition}\label{def:inner-product}
设 $\alpha=(a_1,\cdots,a_n)$、$\beta=(b_1,\cdots,b_n)\in\R^n$，称数
\[
  (\alpha,\beta)=a_1b_1+a_2b_2+\cdots+a_nb_n=\alpha\trans\beta
\]
为 $\alpha$ 与 $\beta$ 的\textbf{内积}。称 $\|\alpha\|=\sqrt{(\alpha,\alpha)}
=\sqrt{a_1^2+\cdots+a_n^2}$ 为 $\alpha$ 的\textbf{长度}（范数）；长度为 $1$ 的
向量称为\textbf{单位向量}。非零向量 $\alpha$ 的单位化向量为
$\alpha_0=\alpha/\|\alpha\|$。
\end{definition}

内积满足对称性 $(\alpha,\beta)=(\beta,\alpha)$、对两个变元的线性性，以及
$(\alpha,\alpha)\ge0$ 且等号成立当且仅当 $\alpha=0$。

\begin{theorem}[柯西--施瓦茨不等式]\label{thm:schwarz}
$|(\alpha,\beta)|\le\|\alpha\|\,\|\beta\|$，等号成立当且仅当 $\alpha,\beta$
线性相关。
\end{theorem}
\begin{proof}
若 $\beta=0$ 显然。设 $\beta\neq0$，对任意实数 $t$，
$0\le(\alpha-t\beta,\alpha-t\beta)=\|\alpha\|^2-2t(\alpha,\beta)+t^2\|\beta\|^2$。
右端是 $t$ 的二次函数且恒非负，故判别式 $\le0$：
$(\alpha,\beta)^2\le\|\alpha\|^2\|\beta\|^2$。
\end{proof}

由柯西--施瓦茨不等式，可以定义两向量夹角 $\theta$ 的余弦
$\cos\theta=(\alpha,\beta)/(\|\alpha\|\|\beta\|)$。当 $(\alpha,\beta)=0$ 时
称 $\alpha$ 与 $\beta$ \textbf{正交}（垂直）。一组两两正交的非零向量称为
\textbf{正交向量组}；若其中每个还是单位向量，称为\textbf{标准正交组}。

\begin{theorem}\label{thm:orthogonal-independent}
正交向量组必线性无关。
\end{theorem}
\begin{proof}
设 $\beta_1,\cdots,\beta_m$ 两两正交且 $k_1\beta_1+\cdots+k_m\beta_m=0$。两边
与 $\beta_j$ 作内积，利用两两正交性得 $k_j\|\beta_j\|^2=0$；$\beta_j\neq0$，
故 $k_j=0$（$j=1,\cdots,m$）。
\end{proof}

\begin{qa}{怎样把一组无关向量“改造”成正交组？减掉的到底是什么？}
答：用\textbf{施密特正交化}。设 $\alpha_1,\cdots,\alpha_m$ 线性无关，令
\[
  \beta_1=\alpha_1,\qquad
  \beta_j=\alpha_j-\sum_{i<j}\frac{(\alpha_j,\beta_i)}{\|\beta_i\|^2}\,\beta_i
  \quad(j=2,\cdots,m).
\]
减掉的第 $i$ 项正是 $\alpha_j$ 在已有正交方向 $\beta_i$ 上的\textbf{投影}
（图~\ref{fig:schmidt}）：$\dfrac{(\alpha_j,\beta_i)}{\|\beta_i\|^2}\beta_i$
是 $\beta_i$ 方向上“与 $\alpha_j$ 内积相同”的那个向量。减完全部投影后，剩余
分量与每个 $\beta_i$ 的内积都是 $0$——垂直是“减掉投影”的直接后果：
\[
  (\beta_j,\beta_i)
  =(\alpha_j,\beta_i)-\frac{(\alpha_j,\beta_i)}{\|\beta_i\|^2}(\beta_i,\beta_i)=0
  \quad(i<j).
\]
最后逐个单位化，即得与原向量组\textbf{张成相同}的标准正交组（每一步只做减法
与数乘，不改变张成空间，故“改造”是无损的）。
\end{qa}

\begin{figure}[htbp]
\centering
\begin{tikzpicture}[>=Stealth, scale=1.15]
  \draw[->, gray!70] (-0.9,0) -- (2.6,0);
  \draw[->, gray!70] (0,-0.5) -- (0,2.3);
  \draw[very thick, blue!60!black, ->] (0,0) -- (1.8,0.6)
        node[below right] {$\beta_1=\alpha_1$};
  \draw[thick, gray!70!black, ->] (0,0) -- (1.2,1.8) node[above] {$\alpha_2$};
  \draw[dashed, blue!60!black, ->] (0,0) -- (1.62,0.54)
        node[below, xshift=2pt] {\small 投影};
  \draw[dashed, gray] (1.2,1.8) -- (1.62,0.54);
  \draw[very thick, red!70!black, ->] (0,0) -- (-0.42,1.26)
        node[left] {$\beta_2=\alpha_2-$ 投影};
  \draw[dashed, red!70!black] (1.62,0.54) -- (1.2,1.8);
  \draw[gray!70] (1.62,0.54) ++(18:0.22) arc[start angle=18, end angle=90, radius=0.22];
\end{tikzpicture}
\caption{施密特正交化：从 $\alpha_2$ 中减去它在 $\beta_1$ 上的投影，剩余分量
$\beta_2$ 与 $\beta_1$ 垂直（图中直角标记处）。}
\label{fig:schmidt}
\end{figure}

\begin{example}\label{ex:schmidt}
把 $\alpha_1=(1,1,0)$、$\alpha_2=(1,0,1)$、$\alpha_3=(0,1,1)$ 正交化并
单位化。
\end{example}
\begin{proof}[解]
\[
  \beta_1=(1,1,0);
  \qquad
  \beta_2=\alpha_2-\tfrac12\beta_1=(\tfrac12,-\tfrac12,1);
  \]
  \[
  \beta_3=\alpha_3-\tfrac12\beta_1-\tfrac13\beta_2
         =\alpha_3-\tfrac12(1,1,0)-\tfrac13(\tfrac12,-\tfrac12,1)
         =(-\tfrac23,\tfrac23,\tfrac23).
\]
去掉正倍数，取 $\beta_1'=(1,1,0)$，$\beta_2'=(1,-1,2)$，$\beta_3'=(-1,1,1)$，
两两正交（可逐一验证内积为零），单位化得
\[
  \gamma_1=\tfrac{1}{\sqrt2}(1,1,0),\quad
  \gamma_2=\tfrac{1}{\sqrt6}(1,-1,2),\quad
  \gamma_3=\tfrac{1}{\sqrt3}(-1,1,1).
\]
\end{proof}

\begin{definition}\label{def:orthogonal-matrix}
满足 $A\trans A=E$ 的 $n$ 阶方阵 $A$ 称为\textbf{正交矩阵}。
\end{definition}

\begin{theorem}\label{thm:orthogonal-matrix}
$A$ 为正交矩阵 $\iff$ $A$ 的列向量组是标准正交组。正交矩阵可逆且
$A^{-1}=A\trans$，$|A|=\pm1$；两个正交矩阵之积仍正交。正交变换保持内积、
长度与夹角不变。
\end{theorem}
\begin{proof}
$A\trans A=E$ 按元素写开即“第 $i$ 列与第 $j$ 列的内积为 $\delta_{ij}$”。
取行列式得 $|A|^2=1$。保持内积：
$(Ax,Ay)=(Ax)\trans(Ay)=x\trans A\trans A\,y=(x,y)$，长度与夹角由内积决定。
\end{proof}

\section{特征值与特征向量}\label{sec:eigen}

\begin{definition}\label{def:eigen}
设 $A$ 为 $n$ 阶方阵。若存在数 $\lambda$ 与\textbf{非零}向量 $\xi$，使
\[
  A\xi=\lambda\xi,
\]
则称 $\lambda$ 是 $A$ 的一个\textbf{特征值}，$\xi$ 是属于 $\lambda$ 的
\textbf{特征向量}。
\end{definition}

\begin{insight}[本质：特征向量 = 变换“不动”的方向]
把 $A$ 看作空间的变换：一般的向量被移动后\textbf{连方向都变了}；特征向量是
幸运者——它只被伸缩，不偏转，特征值 $\lambda$ 就是伸缩的倍数（$\lambda<0$
表示反向）。于是“找特征向量”就是问：\textbf{这个变换有哪些天生不动的方向，
各伸缩多少倍？}（图~\ref{fig:eig-geo}）。若能找到 $n$ 个无关的不动方向，变换
在这组“主轴”下就只是一组独立的伸缩——这正是“对角化”的全部含义。
\end{insight}

\begin{figure}[htbp]
\centering
\begin{tikzpicture}[>=Stealth, scale=0.9]
  \begin{scope}
    \draw[->, gray!70] (-0.4,0) -- (2.6,0);
    \draw[->, gray!70] (0,-0.4) -- (0,2.6);
    \draw[thick, gray!70!black, ->] (0,0) -- (0.8,1.4) node[above left] {$v$};
    \draw[very thick, red!70!black, ->] (0,0) -- (1.6,1.4) node[right] {$Av$};
    \node[below] at (1.1,-0.7) {一般向量：方向被改变};
  \end{scope}
  \draw[->, very thick, gray] (3.6,1.3) -- (4.9,1.3)
        node[midway, above] {$A=\operatorname{diag}(2,1)$};
  \begin{scope}[xshift=6.0cm]
    \draw[->, gray!70] (-0.4,0) -- (2.9,0);
    \draw[->, gray!70] (0,-0.4) -- (0,2.6);
    \draw[thick, blue!60!black, ->] (0,0) -- (1.3,0);
    \draw[very thick, blue!60!black, ->] (0,0) -- (2.6,0)
        node[below] {$\xi$（$\lambda=2$）};
    \draw[very thick, red!70!black, ->] (0,0) -- (0,1.6)
        node[left] {$\eta$（$\lambda=1$）};
    \node[below] at (1.2,-0.8) {特征向量：只伸缩，不偏转};
  \end{scope}
\end{tikzpicture}
\caption{伸缩变换 $\operatorname{diag}(2,1)$：一般向量（左）方向改变；坐标轴
方向（右）天生不动，$x$ 轴方向拉伸 $2$ 倍，$y$ 轴方向保持 $1$ 倍。}
\label{fig:eig-geo}
\end{figure}

\begin{theorem}\label{thm:eigen-criterion}
$\lambda$ 是 $A$ 的特征值 $\iff$ $|A-\lambda E|=0$；$\xi$ 是属于 $\lambda$ 的
特征向量 $\iff$ $\xi$ 是齐次方程组 $(A-\lambda E)x=0$ 的非零解。
\end{theorem}
\begin{proof}
$A\xi=\lambda\xi$ 即 $(A-\lambda E)\xi=0$。它有非零解当且仅当系数矩阵降秩，
即 $|A-\lambda E|=0$（第四章定理~\ref{thm:fundamental} 的特例）。
\end{proof}

称 $f(\lambda)=|A-\lambda E|$ 为 $A$ 的\textbf{特征多项式}。求特征值的步骤：
解特征方程 $f(\lambda)=0$；对每个特征值解 $(A-\lambda E)x=0$，其基础解系给
出全部特征向量。

\begin{example}\label{ex:eigen-2}
求 $A=\begin{pmatrix} 3 & -1 \\ -1 & 3 \end{pmatrix}$ 的特征值与特征向量。
\end{example}
\begin{proof}[解]
\[
  |A-\lambda E|=\begin{vmatrix} 3-\lambda & -1 \\ -1 & 3-\lambda \end{vmatrix}
  =(3-\lambda)^2-1=(2-\lambda)(4-\lambda),
\]
故 $\lambda_1=2$，$\lambda_2=4$。对 $\lambda_1=2$：解 $(A-2E)x=0$，即
$\begin{pmatrix} 1 & -1 \\ -1 & 1 \end{pmatrix}x=0$，得 $\xi_1=(1,1)\trans$
（验证：$A\xi_1=(2,2)\trans=2\xi_1$）。对 $\lambda_2=4$：得
$\xi_2=(1,-1)\trans$（$A\xi_2=(4,-4)\trans=4\xi_2$）。
\end{proof}

\begin{theorem}\label{thm:eigen-sum-product}
设 $n$ 阶方阵 $A$ 的全部特征值为 $\lambda_1,\cdots,\lambda_n$（按重数计），
则
\[
  \lambda_1+\cdots+\lambda_n=\operatorname{tr}(A)=a_{11}+\cdots+a_{nn},
  \qquad
  \lambda_1\lambda_2\cdots\lambda_n=|A|.
\]
\end{theorem}

在特征多项式的展开中，$\lambda^{n-1}$ 的系数由主对角线乘积这一项唯一贡献，
故等于 $-\operatorname{tr}(A)$；令 $\lambda=0$ 得常数项 $f(0)=|A|$；再由根与
系数的关系即得（详细证明从略）。该定理是很好的快速校验：例~\ref{ex:eigen-2}
中 $2+4=6=\operatorname{tr}(A)$，$2\times4=8=|A|$，两条都吻合。

\begin{theorem}\label{thm:distinct-independent}
属于\textbf{互不相同}特征值的特征向量线性无关。
\end{theorem}
\begin{proof}
设 $\lambda_1,\cdots,\lambda_m$ 互不相同，$\xi_i$ 属于 $\lambda_i$，且
$c_1\xi_1+\cdots+c_m\xi_m=0$。对等式两边施加乘积
$(A-\lambda_2E)(A-\lambda_3E)\cdots(A-\lambda_mE)$。因
$(A-\lambda_jE)\xi_j=0$，第 $j\ ( \ge2)$ 项全部被“杀死”；第 1 项变为
\[
  c_1\prod_{j=2}^{m}(\lambda_1-\lambda_j)\,\xi_1=0 .
\]
$\xi_1\neq0$ 且各因子非零，故 $c_1=0$。同理（依次换序施加）可得每个
$c_i=0$。
\end{proof}

\begin{qa}{特征向量为什么必须非零？$\lambda=0$ 又说明什么？}
答：零向量对\textbf{任何} $\lambda$ 都满足 $A\cdot0=\lambda\cdot0$，若允许
它当特征向量，“特征”二字就毫无信息量，故定义明确要求 $\xi\neq0$。$\lambda=0$
成为特征值 $\iff|A-0\cdot E|=|A|=0$，即 $A$ 奇异；此时特征向量正是齐次方程组
$Ax=0$ 的非零解——“被变换压成零的方向”。这是第二章“等价网”的又一条边：
\textbf{有零特征值 $\iff$ 不可逆 $\iff$ 秩亏 $\iff$ 列相关 $\iff$ 压扁}。
\end{qa}

\section{相似矩阵与矩阵的对角化}\label{sec:similar}

\begin{definition}\label{def:similar}
设 $A$、$B$ 为 $n$ 阶方阵。若存在可逆矩阵 $P$ 使 $P^{-1}AP=B$，则称 $B$ 与
$A$ \textbf{相似}，记作 $A\sim B$。若 $B$ 是对角矩阵 $\Lambda$，则称 $A$
\textbf{可对角化}。
\end{definition}

相似满足反身性、对称性、传递性；相似的矩阵同秩、同行列式、同迹、同特征
多项式（从而同特征值）。同特征多项式的验证：
\[
  |B-\lambda E|=|P^{-1}AP-\lambda E|
  =|P^{-1}(A-\lambda E)P|=|P^{-1}|\,|A-\lambda E|\,|P|=|A-\lambda E|.
\]
注意反之不成立：相似比“同特征多项式”苛刻得多。

\begin{insight}[本质：相似 = 同一个变换在不同基下的两张照片]
把 $A$ 看作一个变换，取一组新基排成矩阵 $P$，则同一个变换在新基下的“照片”
恰是 $P^{-1}AP$。相似的矩阵不是两个变换，而是\textbf{一个变换的两个机位}。
对角化就是挑一个最好的机位——让变换看起来最简单：沿 $n$ 个互相独立的方向
各自伸缩 $\lambda_i$，别无杂动作。这个“最好的机位”存在当且仅当 $A$ 有
$n$ 个线性无关的特征向量（下面的定理）。
\end{insight}

\begin{theorem}\label{thm:diagonalizable}
$n$ 阶方阵 $A$ 可对角化 $\iff$ $A$ 有 $n$ 个线性无关的特征向量。此时令
$P=(\xi_1,\cdots,\xi_n)$（列为特征向量），则
\[
  P^{-1}AP=\Lambda=\operatorname{diag}(\lambda_1,\cdots,\lambda_n),
\]
且 $\Lambda$ 的第 $i$ 个对角元属于 $P$ 的第 $i$ 列。特别地，若 $A$ 有 $n$ 个
互不相同的特征值，则 $A$ 必可对角化（逆命题不成立，如 $E$ 本身）。
\end{theorem}
\begin{proof}
把“$A\xi_i=\lambda_i\xi_i$（$i=1,\cdots,n$）”拼成矩阵等式：
$AP=P\Lambda$，即 $P^{-1}AP=\Lambda$（$P$ 的列无关保证 $P$ 可逆）。反之，
由 $P^{-1}AP=\Lambda$ 得 $AP=P\Lambda$，逐列读出 $A\xi_i=\lambda_i\xi_i$，
且 $P$ 可逆保证这 $n$ 个特征向量无关。“互不相同”情形由定理
~\ref{thm:distinct-independent} 立得。
\end{proof}

\begin{qa}{$P$ 的列必须放什么？$\Lambda$ 的对角元顺序可以随便吗？}
答：$P$ 的第 $i$ 列\textbf{必须}是属于 $\lambda_i$ 的特征向量——由证明中的
矩阵等式 $AP=P\Lambda$ 逐列对照即知。因此 $\Lambda$ 的对角元与 $P$ 的列
\textbf{成对绑定}：交换 $\Lambda$ 中两个特征值的位置，必须同时交换 $P$ 中
对应的两个列。单独换其中一边，等式立刻失效。
\end{qa}

\begin{example}\label{ex:diagonalize}
判断 $A=\begin{pmatrix} 4 & 0 & 0 \\ 0 & 3 & 1 \\ 0 & 1 & 3 \end{pmatrix}$
能否对角化；若能，求 $P$ 与 $\Lambda$，并计算 $A^k$。
\end{example}
\begin{proof}[解]
\[
  |A-\lambda E|=(4-\lambda)\bigl[(3-\lambda)^2-1\bigr]
  =(4-\lambda)^2(2-\lambda),
\]
$\lambda_1=\lambda_2=4$，$\lambda_3=2$。对 $\lambda=4$：解 $(A-4E)x=0$，
得 $x_2=x_3$，基础解系 $\xi_1=(1,0,0)\trans$，$\xi_2=(0,1,1)\trans$；对
$\lambda=2$：得 $\xi_3=(0,1,-1)\trans$。三个特征向量线性无关（其中
$\xi_1,\xi_2$ 属于 $\lambda=4$，$\xi_3$ 属于 $\lambda=2$，跨特征值由定理
~\ref{thm:distinct-independent} 保证无关，同特征值内可直接验证无关）。故
\[
  P=\begin{pmatrix} 1 & 0 & 0 \\ 0 & 1 & 1 \\ 0 & 1 & -1 \end{pmatrix},
  \qquad
  P^{-1}AP=\begin{pmatrix} 4 & & \\ & 4 & \\ & & 2 \end{pmatrix}=\Lambda .
\]
由 $A=P\Lambda P^{-1}$ 得 $A^k=P\Lambda^kP^{-1}$，且
$P^{-1}=\begin{pmatrix} 1 & 0 & 0 \\ 0 & \tfrac12 & \tfrac12 \\
0 & \tfrac12 & -\tfrac12 \end{pmatrix}$，于是
\[
  A^k=\begin{pmatrix}
    4^k & 0 & 0 \\
    0 & \dfrac{4^k+2^k}{2} & \dfrac{4^k-2^k}{2} \\
    0 & \dfrac{4^k-2^k}{2} & \dfrac{4^k+2^k}{2}
  \end{pmatrix}.
\]
（取 $k=1$ 可核对与 $A$ 一致。）
\end{proof}

\section{实对称矩阵的正交对角化}\label{sec:symmetric}

一般方阵未必可对角化，但实对称矩阵是一个例外，而且可以要求变换矩阵是
\textbf{正交}矩阵。

\begin{theorem}\label{thm:symmetric-real}
实对称矩阵的特征值都是实数。
\end{theorem}
\begin{proof}
设 $A\trans=A$（实矩阵），$A\xi=\lambda\xi$，$\xi\neq0$（暂允许复向量）。
取共轭转置：$\bar{\xi}\trans A\trans=\bar{\lambda}\bar{\xi}\trans$；由 $A$
实且对称，即 $\bar{\xi}\trans A=\bar{\lambda}\bar{\xi}\trans$。于是
\[
  \bar{\xi}\trans A\xi=\lambda\,\bar{\xi}\trans\xi,
  \qquad
  \bar{\xi}\trans A\xi=(\bar{\xi}\trans A)\xi=\bar{\lambda}\,\bar{\xi}\trans\xi,
\]
两式相减得 $(\lambda-\bar{\lambda})\,\bar{\xi}\trans\xi=0$。而
$\bar{\xi}\trans\xi=\sum|x_i|^2>0$，故 $\lambda=\bar{\lambda}$ 为实数。
\end{proof}

\begin{theorem}\label{thm:symmetric-orthogonal}
实对称矩阵属于不同特征值的特征向量彼此正交。
\end{theorem}
\begin{proof}
设 $A\xi_1=\lambda_1\xi_1$，$A\xi_2=\lambda_2\xi_2$，$\lambda_1\neq\lambda_2$。
则
\[
  \lambda_1(\xi_1,\xi_2)=\xi_2\trans A\xi_1
  =\xi_2\trans A\trans\xi_1=(A\xi_2)\trans\xi_1=\lambda_2(\xi_1,\xi_2),
\]
故 $(\lambda_1-\lambda_2)(\xi_1,\xi_2)=0$，即 $(\xi_1,\xi_2)=0$。
\end{proof}

\begin{theorem}[谱定理，证明从略]\label{thm:spectral}
设 $A$ 为 $n$ 阶实对称矩阵，则存在正交矩阵 $Q$，使
\[
  Q\trans A\,Q=Q^{-1}AQ=\Lambda=\operatorname{diag}(\lambda_1,\cdots,\lambda_n),
\]
其中 $\lambda_i$ 为 $A$ 的全部特征值（按重数计）。
\end{theorem}

操作步骤：求出全部特征值；对每个特征值求 $(A-\lambda E)x=0$ 的基础解系
（同一特征值内的解向量用施密特正交化——跨特征值的自动正交）；全部单位化
后排成 $Q$ 的列。

\begin{insight}[本质：实对称 = 沿互相垂直的主轴纯拉伸]
谱定理的几何内容：对称变换没有任何“旋转、剪切”的杂动作，它只是沿 $n$ 条
\textbf{互相垂直的主轴}各自伸缩。对比一下损失在哪：旋转矩阵
$\begin{pmatrix} 0 & -1 \\ 1 & 0 \end{pmatrix}$ 把每个方向都转走 $90^\circ$，
一个实特征方向都没有——所以“不可对角化/没有实特征值”并不可耻，只是这类
变换的“不动方向”不在这个实平面上。
\end{insight}

\begin{example}\label{ex:spectral}
设 $A=\begin{pmatrix} 0 & 1 & 1 \\ 1 & 0 & 1 \\ 1 & 1 & 0 \end{pmatrix}$，
求正交矩阵 $Q$ 使 $Q\trans AQ$ 为对角矩阵。
\end{example}
\begin{proof}[解]
$|A-\lambda E|=-(\lambda-2)(\lambda+1)^2$（各行加到第一行提出
$(2-\lambda)$ 再化简），$\lambda_1=2$，$\lambda_{2,3}=-1$（二重）。对
$\lambda=2$：解 $(A-2E)x=0$ 得 $\xi_1=(1,1,1)\trans$。对 $\lambda=-1$：
解 $(A+E)x=0$，即 $x_1+x_2+x_3=0$，取 $\eta_2=(1,-1,0)\trans$，
$\eta_3=(1,0,-1)\trans$，施密特正交化：
\[
  \beta_2=\eta_2=(1,-1,0)\trans,\qquad
  \beta_3=\eta_3-\tfrac12\beta_2=(\tfrac12,\tfrac12,-1)\trans .
\]
单位化：
\[
  Q=\left(
  \tfrac{1}{\sqrt3}\xi_1,\ \tfrac{1}{\sqrt2}\beta_2,\ \tfrac{1}{\sqrt6}(1,1,-2)\trans
  \right)
  =\begin{pmatrix}
    \tfrac{1}{\sqrt3} & \tfrac{1}{\sqrt2} & \tfrac{1}{\sqrt6} \\
    \tfrac{1}{\sqrt3} & -\tfrac{1}{\sqrt2} & \tfrac{1}{\sqrt6} \\
    \tfrac{1}{\sqrt3} & 0 & -\tfrac{2}{\sqrt6}
  \end{pmatrix},
\]
则 $Q\trans AQ=\operatorname{diag}(2,-1,-1)$（$Q$ 的列标准正交可逐一验证）。
\end{proof}

\chapterexercises{五}

\begin{enumerate}[label=\arabic*., itemsep=2pt]
\item 求 $\begin{pmatrix} 2 & 1 \\ 1 & 2 \end{pmatrix}$ 的特征值与特征向量。
\item 把 $\alpha_1=(1,1,1)$、$\alpha_2=(1,2,3)$ 施密特正交化并单位化。
\item 判断下列矩阵是否正交矩阵：
  $\text{(1)}\ \dfrac{1}{\sqrt2}\begin{pmatrix} 1 & -1 \\ 1 & 1 \end{pmatrix}$；
  \quad $\text{(2)}\ \begin{pmatrix} 1 & 1 \\ 0 & 1 \end{pmatrix}$。
\item 求 $\begin{pmatrix} 2 & 0 & 0 \\ 0 & 3 & 1 \\ 0 & 1 & 3 \end{pmatrix}$
  的特征值与特征向量，并判断它是否可对角化。
\item 利用第 1 题的结果计算 $\begin{pmatrix} 2 & 1 \\ 1 & 2 \end{pmatrix}^{10}$。
\item 求正交矩阵 $Q$ 使 $Q\trans\begin{pmatrix} 0 & 1 \\ 1 & 0 \end{pmatrix}Q$
  为对角矩阵。
\item 证明：$A$ 与 $A\trans$ 有相同的特征多项式。
\item 设 $A^2=A$（幂等矩阵），证明 $A$ 的特征值只能是 $0$ 或 $1$。
\item 证明：旋转矩阵 $\begin{pmatrix} 0 & -1 \\ 1 & 0 \end{pmatrix}$ 在实数
  范围内没有特征向量，并给出几何解释。
\item 证明：正交矩阵的乘积仍是正交矩阵，且正交矩阵的行列式为 $\pm1$。
\end{enumerate}

% \include{chapters/ch6-quadratic}
% \include{chapters/appendix-answers}

\end{document}
